% ============================================================
% PT Article M4 — version française
% Miroir FR de m4_structures.tex (canonique EN)
% ============================================================

\documentclass[12pt,a4paper]{article}
\IfFileExists{pt-common-fr.sty}{\usepackage{pt-common-fr}}{\usepackage{../shared/pt-common-fr}}
\usepackage[margin=2.5cm]{geometry}
\usepackage{setspace}
\usepackage{tikz-3dplot}
\onehalfspacing

\title{\textbf{La théorie de la persistance~IV~:\\
structures mathématiques étendues}}
\author{Yan Senez\\\texttt{yan.senez@protonmail.com}}
\date{Avril 2026}

\begin{document}
\maketitle

\begin{abstract}
Ceci est le quatrième de cinq articles présentant les fondations
mathématiques de la théorie de la persistance (PT). Construisant sur
la dynamique du crible et les lois de
conservation~\cite{companion_M1}, la géométrie de l'information et
l'holonomie~\cite{companion_M2}, et la théorie de la convergence et
le point fixe~\cite{companion_M3}, nous développons les structures
mathématiques étendues que le crible d'Ératosthène porte au-delà des
théorèmes arithmétiques fondamentaux.

\textbf{Algèbre du crible et transformées.}
Le produit du crible $*_T$ définit une nouvelle algèbre de fonctions
arithmétiques~: non associative, non commutative et sans
élément neutre~--- distincte de la convolution de Dirichlet, des
algèbres de Hecke et de $\ell^1(\mathbb{N})$. La transformée de
persistance $P_K$ projette sur la base propre de la matrice de
transfert $T_3$~; la métrique PT $d_{\rm PT}$ est quasi-orthogonale aux
métriques archimédienne et $p$-adique~; la catégorie
$\mathrm{Cat}(\mathrm{Sieve})$ admet quatre foncteurs cohérents
convergeant vers la structure asymptotique du crible. Trois nouvelles
transformées~--- tensorielle $\mathcal{T}$ (corrélations
inter-modulaires dans un tenseur $3{\times}5{\times}7$), holonomique
$\mathcal{H}$ (produit eulérien généralisé avec métrique de Fisher)
et décohérence $\mathcal{D}$ (monotonie de l'entropie~: la seconde
loi du crible, théorème~H)~--- étendent la transformée de persistance
classique. Le code de résidus de sauts atteint la distance $d = n - 1$
(asymptotiquement MDS)~; prédiction et correction d'erreur coïncident.

\textbf{Mécanique complexe.}
L'observable réelle $\sinsp{p}$ se relève en une variable complexe
$w_p$ sur un cercle de rayon $s = 1/2$, avec
$|w_p|^2 = \mathrm{Re}(w_p)$. La force holomorphe $F(w) = i/w - 2i$,
l'intégrabilité de Liouville sur $\mathbb{T}^n$, une relation
d'incertitude exacte $|p_w|\cdot|w| = 1$, et les corrections
radiatives QED $3\alpha/(4\pi)$ et $16\pi^2$ émergent de la géométrie
du cercle.

\textbf{Mathématiques prédictives.}
Le cadre MP quantifie les cascades de contraintes~: chaque premier
ajoute exactement un degré de liberté (théorème~L1, inconditionnel)~;
la dimension de couche suit $\dim(d) = P(d{+}1) - (2d{+}1)$
(13~strates vérifiées)~; le spectre d'activation est
$\Sigma_k = \binom{4}{k}$ pour $k \geq \mathrm{AT}$~; et la formule de
seuil d'activation en trois régimes (écho, construction, asymptotique)
prouve $\mathrm{AT} = 1$ pour tout $p \geq 41$.

\textbf{Extensions de convergence.}
La course arithmético-géométrique fournit cinq routes alternatives vers
T4. Le graphe de transition modulo~3 $K_3$ est un expanseur de
Ramanujan optimal dont la fonction zêta d'Ihara a tous ses zéros non
triviaux sur le cercle critique. Le groupe diédral
$D_4 = \langle T_3, J \rangle$ organise la structure spectrale via
l'entrelacement~$J$.
\end{abstract}

\tableofcontents
\newpage

\section{Introduction}
\label{sec:intro}

Cet article est le quatrième volet (M4) d'une série en cinq parties
développant les fondations mathématiques de la théorie de la
persistance (PT). La série s'organise comme suit~:
\begin{description}[nosep]
\item[M1~\cite{companion_M1}~:] Le crible comme champ dynamique,
  unicité d'Ératosthène, lois de conservation, et théorème
  fondamental des sauts.
\item[M2~\cite{companion_M2}~:] Géométrie de l'information, holonomie
  et échelle d'angle interne.
\item[M3~\cite{companion_M3}~:] Convergence, récurrences exactes, et
  point fixe $\mustar = 15$.
\item[M4 (cet article)~:] Structures mathématiques étendues~: algèbre
  du crible, mécanique complexe, mathématiques prédictives, et
  extensions de convergence.
\item[M5~\cite{companion_M5}~:] Le pont de l'arithmétique vers la
  physique.
\end{description}

\noindent
Ces cinq articles fournissent ensemble une fondation mathématique
auto-contenue de la théorie de la persistance, destinée à être lue
dans l'ordre M1\,$\to$\,M5.

Les articles compagnons~\cite{companion_M1,companion_M2,companion_M3}
ont établi la théorie arithmétique de base~: la dynamique du crible
et les lois de conservation (M1), la géométrie de l'information
canonique et les angles d'holonomie
$\sinsp{p} = \delta_p(2 - \delta_p)$ (M2), et la convergence vers
l'unique point fixe $\mustar = 15$ avec premiers actifs
$\{3, 5, 7\}$ (M3). Cet article développe l'infrastructure
mathématique étendue que le crible porte au-delà de ces théorèmes de
base.

Le matériau est organisé en quatre sections, correspondant à quatre
domaines mathématiques distincts.

\textbf{Section~\ref{sec:math_structures}~: Structures mathématiques.}
Le produit du crible $*_T$ pondère la convolution de Dirichlet par la
matrice de transfert $T_3$, produisant une nouvelle structure
algébrique. La transformée de persistance $P_K$, la métrique PT
$d_{\rm PT}$, la catégorie du crible et trois nouvelles transformées
(tensorielle, holonomique, décohérence) révèlent la profondeur du
contenu mathématique du crible. Le code des résidus de sauts atteint
la distance de Hamming $d = n - 1$ (théorème, asymptotiquement
MDS)~: 32~scripts compagnons valident le catalogue complet.

\textbf{Section~\ref{sec:complex}~: Mécanique complexe.}
Le relèvement de l'observable réelle $\sinsp{p}$ à la variable
complexe $w_p = (1 - e^{2i\theta_p})/2$ met au jour une structure
cachée~: tous les $w_p$ se trouvent sur un cercle de rayon $s = 1/2$
(théorème du cercle), la force $F(w) = i/w - 2i$ est holomorphe, le
moment angulaire vaut $\sinsp{p}$ (charge de Noether), et le système
est intégrable au sens de Liouville sur $\mathbb{T}^n$. Les
corrections radiatives QED se décomposent entièrement en quantités
PT-natives.

\textbf{Section~\ref{sec:PM}~: Mathématiques prédictives.}
Le cadre MP prédit ce qui devient structurellement inévitable
lorsque des contraintes sont empilées sur un champ brut. Sept
théorèmes à partir d'une unique fondation algébrique (algèbre des
indicateurs + TCR)~: L1 (incrément unité), la formule de dimension,
le spectre d'activation, la formule AT à trois régimes, l'identité de
partition C1, la saturation fantôme C2 et le corollaire MP--PT
($\mathrm{AT} = 1$ pour tout $p \geq 41$). Trois instanciations
(crible, codes correcteurs d'erreurs, nombres chanceux) affinent la
classification universel/spécifique.

\textbf{Section~\ref{sec:convergence_ext}~: Extensions de
convergence.}
La course arithmético-géométrique, la théorie spectrale des graphes
du crible (expanseur de Ramanujan optimal, fonction zêta d'Ihara) et
le groupe de symétrie $D_4$ approfondissent et valident le théorème
de convergence T4 de~\cite{companion_M3}.

Tous les résultats découlent de $s = 1/2$ et de la structure de crible
établie dans~\cite{companion_M1,companion_M2,companion_M3}. Aucune
physique n'est invoquée~; aucun paramètre n'est ajusté.

\section{Structures mathématiques du crible}
\label{sec:math_structures}

\epigraph{Le but du calcul est la compréhension, pas les
nombres.}{Richard Hamming}

% ---- Status Box ----------------------------------------
\begin{statusbox}
\textbf{OBJECTIF~:}
  Exposer les structures algébriques, métriques, catégorielles,
  quantiques et de théorie des codes que le crible porte au-delà des
  applications physiques des articles
  compagnons~\cite{companion_M3,companion_physics}.

\textbf{ENTRÉES~:}
  $T_3$ (article compagnon~\cite{companion_M1}), TFS (article
  compagnon~\cite{companion_M1}), géométrie de Fisher (article
  compagnon~\cite{companion_M2}), structure TCR (article
  compagnon~\cite{companion_M1}), borne spectrale (article
  compagnon~\cite{companion_M3}).

\textbf{RÉSULTATS PRINCIPAUX~:}
  \begin{itemize}[nosep]
  \item Nouvelle algèbre du crible $*_T$~: non associative, non
    commutative, sans élément neutre (M15).
  \item Transformée de persistance $P_K$~: décomposition spectrale
    exacte (M16).
  \item Métrique PT $d_{\rm PT}$~: quasi-orthogonale à $|\cdot|$ et
    $|\cdot|_p$ (M17).
  \item Nombres PT $\mathbb{Z}_{\rm PT}$~: enrichissement profini de
    $\mathbb{Z}$ (M18).
  \item Catégorie $\mathrm{Cat}(\mathrm{Sieve})$~: quatre foncteurs
    cohérents (M19).
  \item Extension modulaire universelle~: les structures valent pour
    tout premier $q$ (M21).
  \item Crible quantique~: canal CPTP avec décohérence arithmétique
    (M27).
  \item Réseau tensoriel TCR~: MPS de dimension de liaison 3 exacte
    (M28).
  \item Prédiction multi-modulaire~: $+61\%$ par rapport à mod-3 seul
    (M30).
  \item Code quantique TCR $[[3, 5.58, 1]]$~: prédiction = correction
    d'erreur (M31).
  \item Code de distance de sauts~: $d = n - 1$ (THÉORÈME),
    asymptotiquement MDS (M32).
  \item Transformée tensorielle $\mathcal{T}$~: $3{\times}5{\times}7
    = 105$ composantes, le rang classifie les corrélations
    inter-modulaires (M33).
  \item Transformée holonomique $\mathcal{H}$~: $f \mapsto
    \{\sinp{p}(f)\}$, métrique de Fisher, produit eulérien généralisé
    (M34).
  \item Transformée de décohérence $\mathcal{D}$~: $f \mapsto
    \{S_K(f)\}$, théorème~H (monotonie = 2\textsuperscript{e}~loi du
    crible) (M35).
  \end{itemize}

\textbf{STATUT~:}
  \VALTAG~(32 scripts compagnons~; cf.\ scripts compagnons).

\textbf{NON PROUVÉ~:}
  L'extension quantique CSS donne $d_{\rm CSS} \sim 1$ pour $n$~grand
  (la distance classique $d = n-1$ ne se transfère pas pleinement au
  cadre quantique). Les applications cryptographiques sont
  spéculatives.
\end{statusbox}


% ============================================================
\subsection{L'algèbre du crible}
\label{sec:ms_algebra}

\begin{definition}[Produit du crible $*_T$]
\label{def:sieve_product}
Pour des fonctions arithmétiques $f, g\colon \mathbb{N} \to
\mathbb{R}$, on définit~:
\begin{equation}
  (f *_T g)(n) = \sum_{d \mid n} f(d)\, g(n/d)\,
    T_3\bigl(d \bmod 3,\; (n/d) \bmod 3\bigr).
  \label{eq:sieve_product}
\end{equation}
\end{definition}

Le produit du crible pondère chaque paire de diviseurs par l'entrée de
la matrice de transfert reliant leurs classes de résidus. Il hérite
de la structure divisorielle de la convolution de Dirichlet, mais la
module par la dynamique du crible.

\begin{thmbox}{Algèbre du crible}
\label{thm:sieve_algebra}
L'algèbre $(\mathcal{A}, *_T)$ des fonctions arithmétiques sous le
produit du crible possède les propriétés suivantes~:
\begin{enumerate}[nosep]
\item \textbf{Non associative}~: $(f *_T g) *_T h \neq f *_T (g *_T h)$
  en général (vérifié par contre-exemple explicite pour $n = 6$).
\item \textbf{Non commutative}~: $f *_T g \neq g *_T f$ en général
  (asymétrie de $T_3$~: $T_{01} \neq T_{10}$).
\item \textbf{Sans élément neutre}~: il n'existe pas de fonction $e$
  telle que $e *_T f = f *_T e = f$ pour tout~$f$.
\item \textbf{Distincte de toutes les algèbres connues}~: $*_T$
  diffère de la convolution de Dirichlet, des algèbres de Hecke et de
  $\ell^1(\mathbb{N})$ par sa non-associativité et sa pondération par
  matrice de transition.
\end{enumerate}
Script~: \texttt{ch\_math\_structures/test\_sieve\_algebra.py}
(scripts compagnons).
\end{thmbox}


% ============================================================
\subsection{La transformée de persistance}
\label{sec:ms_transform}

\begin{definition}[Transformée de persistance $P_K$]
\label{def:persistence_transform}
Pour une fonction arithmétique $f$ et un niveau de crible $K$, on
définit $P_K(f)$ comme la projection spectrale de la restriction
$f|_{\mathrm{surv}(K)}$ sur la base propre $(v_+, v_-)$ de~$T_3$~:
\begin{equation}
  P_K(f) = \langle f, v_+ \rangle_K \cdot v_+
          + \langle f, v_- \rangle_K \cdot v_-,
  \label{eq:persistence_transform}
\end{equation}
où $\langle \cdot, \cdot \rangle_K$ est le produit scalaire pondéré
par la distribution des survivants au niveau~$K$.
\end{definition}

Propriétés clés~:
\begin{itemize}[nosep]
\item $P_K(\mathbf{1}) = v_+$ (la fonction constante est purement
  stationnaire).
\item $P_K(\chi_3) = v_-$ (le caractère de Legendre est purement
  anti-stationnaire).
\item $P_K$ est linéaire, idempotente ($P_K^2 = P_K$) et exacte
  (inversible sur le sous-espace spectral de dimension~2).
\item $P_K$ est compatible avec $*_T$~:
  $P_K(f *_T g) = P_K(f) *_T P_K(g)$ sur l'espace propre.
\end{itemize}

Script~: \texttt{ch\_math\_structures/test\_persistence\_transform.py}
(scripts compagnons).


% ============================================================
\subsection{La métrique PT et les nombres PT}
\label{sec:ms_metric}

\subsubsection{La métrique PT}

\begin{definition}[Distance PT]
\label{def:pt_metric}
Pour des entiers $m, n$, on définit~:
\begin{equation}
  d_{\rm PT}(m, n) = \sum_{K=2}^{K_{\max}}
    \bigl| \sigma_K(m) - \sigma_K(n) \bigr|,
  \label{eq:pt_metric}
\end{equation}
où $\sigma_K(m) = m \bmod p_K$ est la signature du crible au
niveau~$K$ et $K_{\max}$ est le niveau de crible le plus profond
considéré.
\end{definition}

\begin{thmbox}{Propriétés de $d_{\rm PT}$}
\label{thm:pt_metric}
\begin{enumerate}[nosep]
\item $d_{\rm PT}$ est une métrique (inégalité triangulaire, symétrie,
  $d_{\rm PT}(m,n) = 0 \Leftrightarrow m = n$ dans la plage
  primorielle).
\item $d_{\rm PT}$ est \textbf{quasi-orthogonale} à la fois à la
  métrique archimédienne $|m - n|$ et aux métriques $p$-adiques
  $|m-n|_p$~: corrélation $< 0{,}1$ dans tous les cas.
\item Sous $d_{\rm PT}$, les premiers sont \textbf{dispersés} (deux
  premiers ne sont jamais proches) tandis que les composés sont
  \textbf{regroupés} (partage de motifs de divisibilité). Précision de
  classification~: $96\%$.
\end{enumerate}
Script~: \texttt{ch\_math\_structures/test\_pt\_metric.py} (scripts
compagnons).
\end{thmbox}

\subsubsection{Nombres PT}

\begin{definition}[Enrichissement PT $\mathbb{Z}_{\rm PT}$]
\label{def:pt_numbers}
On définit $\mathbb{Z}_{\rm PT}$ comme l'ensemble des entiers enrichis
par leur trajectoire dans le crible~:
$n \mapsto (n,\, \sigma_2(n),\, \sigma_3(n),\, \ldots)$.
La multiplication est héritée de $\mathbb{Z}$~; l'addition est
«~disruptive~» (elle brise la trajectoire dans le crible).
\end{definition}

Propriétés clés~: $\sigma_K$ est un homomorphisme d'anneaux, donc
$\sigma_K(m+n) = \sigma_K(m) + \sigma_K(n) \bmod p_K$ et
$\sigma_K(mn) = \sigma_K(m) \cdot \sigma_K(n) \bmod p_K$ tiennent
exactement. Cependant, le \emph{statut de survivant} du crible (si
$m$ est $K$-rugueux) n'est pas préservé par l'addition~: deux entiers
$K$-rugueux peuvent se sommer en un composé. $\mathbb{Z}_{\rm PT}$
est un \emph{enrichissement filtré profini} de $\mathbb{Z}$.

Script~: \texttt{ch\_math\_structures/test\_pt\_numbers.py} (scripts
compagnons).


% ============================================================
\subsection{Catégorie du crible}
\label{sec:ms_category}

\begin{definition}[$\mathrm{Cat}(\mathrm{Sieve})$]
\label{def:sieve_category}
La \emph{catégorie du crible} a~:
\begin{itemize}[nosep]
\item Objets~: niveaux du crible $K = 2, 3, 4, \ldots$
\item Morphismes~: applications de transition $\phi_{K \to K+1}$
  induites par le premier suivant $p_{K+1}$.
\end{itemize}
\end{definition}

Quatre foncteurs cohérents relient le crible aux catégories
classiques~:
\begin{enumerate}[nosep]
\item $F_{\rm Grp}\colon \mathrm{Cat}(\mathrm{Sieve}) \to \mathbf{Grp}$~:
  envoie le niveau $K$ sur le groupe de symétrie $G_K$ de la matrice
  de transition.
\item $F_{\rm Vect}\colon \mathrm{Cat}(\mathrm{Sieve}) \to \mathbf{Vect}$~:
  envoie le niveau $K$ sur l'espace propre $V_K$ de~$T_K$.
\item $F_{\rm Top}\colon \mathrm{Cat}(\mathrm{Sieve}) \to \mathbf{Top}$~:
  envoie le niveau $K$ sur le complexe simplicial du graphe de
  transition.
\item $F_{\rm Info}\colon \mathrm{Cat}(\mathrm{Sieve}) \to \mathbf{Info}$~:
  envoie le niveau $K$ sur l'espace de divergence
  $(D_{\rm KL}^{(K)}, g_F^{(K)})$.
\end{enumerate}

Les transformations naturelles entre ces foncteurs encodent la
compatibilité des descriptions algébriques, spectrales, topologiques
et information-théoriques. La limite projective $\varprojlim_K F(K)$
converge pour les quatre foncteurs, donnant la structure asymptotique
du crible.

Script~: \texttt{ch\_math\_structures/test\_sieve\_category.py}
(scripts compagnons).


% ============================================================
\subsection{Extension modulaire universelle}
\label{sec:ms_modular}

Les structures décrites ci-dessus ont été construites pour le crible
modulo~3. Une question naturelle~: dépendent-elles de $q = 3$~?

\begin{thmbox}{Universalité du produit du crible}
\label{thm:universality}
Pour tout premier impair $q$, la matrice de transfert du crible
modulo~$q$ $T_q$ présente~:
\begin{enumerate}[nosep]
\item Un trou spectral~: $|\lambda_2(T_q)| < 1$.
\item Des auto-transitions interdites~: $T_q(r, r) = 0$ pour $r \neq 0$.
\item Une algèbre du crible $*_{T_q}$ avec les mêmes propriétés
  algébriques (non associative, non commutative, sans élément neutre).
\item Une transformée de persistance $P_K^{(q)}$ avec décomposition
  spectrale.
\end{enumerate}
Vérifié pour $q = 3, 5, 7, 11, 13$. Toutes les structures sont
\textbf{universelles}~--- elles dépendent de la primalité de $q$, et
non spécifiquement de $q = 3$.
Script~: \texttt{ch\_math\_structures/test\_modular\_extension.py}
(scripts compagnons).
\end{thmbox}


% ============================================================
\subsection{Transformées PT (M33--M35)}
\label{sec:ms_new_transforms}

La transformée de persistance~$P_K$ (\cref{sec:ms_transform}) projette
sur la base propre de~$T_3$, capturant le contenu spectral modulo~3
en deux coordonnées par~$K$. Trois transformées complémentaires
étendent cette image en capturant les \emph{corrélations
inter-modulaires} (M33), la \emph{géométrie holonomique} (M34) et la
\emph{monotonie entropique}~(M35).


\subsubsection{La transformée tensorielle (M33)}
\label{sec:ms_tensor_transform}

\begin{definition}[Transformée tensorielle de persistance]
\label{def:tensor_transform}
Pour une fonction arithmétique~$f$ et une profondeur de crible~$K$,
soit $P_K^{(q)}(f) \in \mathbb{R}^q$ la projection spectrale de~$f$
sur la base propre de la matrice de transfert~$T_q$ (cf.~M21). La
\emph{transformée tensorielle de persistance} est
\begin{equation}
  \mathcal{T}_K(f)
  = P_K^{(3)}(f) \otimes P_K^{(5)}(f) \otimes P_K^{(7)}(f)
  \;\in\; \mathbb{R}^{3 \times 5 \times 7},
  \label{eq:tensor_transform}
\end{equation}
un tenseur d'ordre trois à $3 \times 5 \times 7 = 105$ composantes.
\end{definition}

La concaténation multi-modulaire (M30, \cref{sec:ms_codes}) empile les
trois vecteurs de projection en un vecteur de dimension~$15$,
abandonnant toute information inter-modulaire. Le produit tensoriel
conserve les \emph{corrélations inter-modulaires}~: comment $f$ se
comporte \emph{simultanément} modulo~$3$, modulo~$5$ et modulo~$7$.

\begin{thmbox}{Classification par rang tensoriel}
\label{thm:tensor_rank}
\begin{enumerate}[nosep]
\item $f = \mathbf{1}$~: $\mathcal{T}_K(\mathbf{1})$ a rang~$1$ pour
  tout~$K$ (produit pur, aucune corrélation inter-modulaire).
\item $f = \chi_3$~: rang $\leq 2$ (couplage faible, structure modulo~3
  découplable).
\item $f = \lambda, \mu$ (Liouville, Möbius)~: rang $2$--$3$
  (couplage arithmétique non trivial entre modules).
\end{enumerate}
Le rang tensoriel mesure la \textbf{complexité} des corrélations
inter-modulaires induites par le crible.
\end{thmbox}

\paragraph{Identité de Plancherel.}
Pour un tenseur de rang~1~:
\begin{equation}
  \|\mathcal{T}_K(f)\|^2
  = \|P_K^{(3)}\|^2 \cdot \|P_K^{(5)}\|^2 \cdot \|P_K^{(7)}\|^2.
  \label{eq:plancherel_tensor}
\end{equation}
Le rapport de couplage $\rho = \|\mathcal{T}\|^2 / \prod_q
\|P^{(q)}\|^2$ vaut exactement~$1$ lorsque les facteurs TCR sont
informationnellement indépendants~; les écarts quantifient les
corrélations induites par le crible.

\paragraph{Décomposition CP.}
La décomposition polyadique canonique
$\mathcal{T}_{jkl} = \sum_{r=1}^{R} a_{r,j}^{(3)} \, a_{r,k}^{(5)} \,
a_{r,l}^{(7)}$
sépare le tenseur en $R$~composantes de rang~1~: le mode $r = 1$ est
la part TCR-factorisable~; les modes $r \geq 2$ sont des termes de
correction inter-modulaires.

Script~: \texttt{ch\_math\_structures/test\_tensor\_transform.py}
(scripts compagnons).


\subsubsection{La transformée holonomique (M34)}
\label{sec:ms_holonomic_transform}

\begin{definition}[Transformée holonomique $\mathcal{H}$]
\label{def:holonomic_transform}
Pour une fonction arithmétique~$f$, une profondeur de crible~$K$ et
un premier actif~$p$ ($p = p_j$, $j \leq K$), on définit le
\emph{déficit généralisé}
\begin{equation}
  \delp{p}(f,K) = 1 - \frac{w_0}{\sum_{c=0}^{p-1} w_c},
  \qquad
  w_c = \frac{n_c}{N}\,\langle |f|^2 \rangle_c,
  \label{eq:holonomic_deficit}
\end{equation}
où $n_c$ compte les sauts dans la classe~$c \bmod p$ et
$\langle |f|^2 \rangle_c$ est la moyenne conditionnelle de~$|f|^2$.
L'\emph{angle holonomique} et son sinus carré sont
\begin{equation}
  \cos\theta_p(f) = 1 - \delp{p}(f),
  \qquad
  \sinp{p}(f) = \delp{p}(f)\,(2 - \delp{p}(f)).
  \label{eq:holonomic_angle}
\end{equation}
La \emph{transformée holonomique} est l'application
$\mathcal{H}(f,K) = \bigl\{\sinp{p}(f,K)\bigr\}_{p\;\text{actif}}$.
\end{definition}

\begin{thmbox}{Propriétés de la transformée holonomique}
\label{thm:holonomic_properties}
\begin{enumerate}[nosep]
\item \textbf{Identité pythagoricienne}~:
  $\sinp{p} + \cos^2\!\theta_p = 1$ pour tout~$p$ (par construction).
\item \textbf{Produit eulérien généralisé}~:
  $\Pi(f,K) = \prod_{p\;\text{actif}} \sinp{p}(f,K)$.
  Pour $f = \mathbf{1}$, c'est l'holonomie brute du crible.
  Pour les poids $\qplus$, $\Pi \to \alphaEM \approx 1/137$.
\item \textbf{Métrique de Fisher}~: la variation infinitésimale
  \begin{equation}
    ds^2 = \sum_p \frac{(d\theta_p)^2}{\sinp{p}}
    \label{eq:fisher_holonomic}
  \end{equation}
  est une métrique riemannienne sur l'espace des fonctions
  arithmétiques, mesurant la courbure informationnelle dans l'espace
  du crible (cf.\ article compagnon~\cite{companion_M2}).
\item \textbf{Stabilité}~: pour $p$ fixé, $\theta_p(f,K)$ converge
  lorsque $K \gg \mathrm{index}(p)$.
\item \textbf{Distance holonomique signée}~: pour des fonctions avec
  $|f|^2 = |g|^2$ (par ex.\ $f = \mathbf{1}$, $g = \chi_3$), le
  spectre fondé sur l'énergie est identique. La distance
  \emph{signée} utilise les moyennes par classe
  $m_c^f = \langle f \rangle_c$~:
  \begin{equation}
    d_{\mathcal{H}}(f,g)
    = \sqrt{\sum_p \sum_{c=0}^{p-1}
      \frac{(m_c^f - m_c^g)^2}{p}}.
    \label{eq:holonomic_distance}
  \end{equation}
\end{enumerate}
Script~: \texttt{ch\_math\_structures/test\_holonomic\_transform.py}
(scripts compagnons).
\end{thmbox}


\begin{remarque}[Lecture circulaire de la transformée holonomique]\mBR
\label{rem:circular_reading}
L'angle $\theta_p$ est une véritable variable circulaire sur $U(1)$,
et non simplement un paramétrage commode d'une quantité réelle. Cela
se voit depuis le relèvement complexe~:
\begin{equation}
  w_p = \frac{1 - e^{2i\theta_p}}{2}
  = \sinp{p}\,e^{i(\pi/2 - \theta_p)},
  \label{eq:complex_lift}
\end{equation}
qui satisfait $(\mathrm{Re}\,w_p - \tfrac{1}{2})^2
+ \mathrm{Im}(w_p)^2 = \tfrac{1}{4}$~: tout $w_p$ se trouve sur le
cercle de centre et rayon $1/2$. La valeur $1/2$ qui apparaît ici est
\emph{le même} $s = 1/2$ que l'occupation stationnaire modulo~3~:
les deux expriment le point d'équilibre de la dynamique de
transitions interdites du crible.

L'identité $\sinp{p} = \mathrm{Re}(1 - e^{2i\theta_p})/2$ fait de la
structure pythagoricienne (item~1 ci-dessus) une projection de la
géométrie circulaire sur la droite réelle, tandis que la métrique de
Fisher \eqref{eq:fisher_holonomic} vit naturellement sur le tore
produit $\prod_p U(1)$. Un développement complet des fondations
circulaires (donnée~$C_0$~: phase, enroulement, holonomie comme
objets primitifs) est reporté à un article compagnon.
\end{remarque}

\subsubsection{La transformée de décohérence (M35)}
\label{sec:ms_decoherence_transform}

\begin{definition}[Transformée de décohérence $\mathcal{D}$]
\label{def:decoherence_transform}
Pour une fonction arithmétique~$f$ et une profondeur de crible~$K$,
on définit~:
\begin{enumerate}[nosep]
\item La \emph{distribution d'énergie par classe}~:
  $p_c(f,K) = \sum_{n:\,\mathrm{classe}(n)=c} |f(n)|^2
   \big/ \sum_n |f(n)|^2$.
\item L'\emph{état quantique} (règle de Born)~:
  $|\psi_K(f)\rangle = \sum_c \sqrt{p_c}\,e^{i\phi_c}\,|c\rangle$,
  où $\phi_c = \arg\langle f \rangle_c$.
\item L'\emph{entropie de décohérence}~:
  $S_K(f) = -\sum_c p_c(f,K)\,\ln p_c(f,K)$.
\end{enumerate}
La \emph{transformée de décohérence} est $\mathcal{D}(f)
= \{S_K(f)\}_{K \geq K_{\min}}$.
\end{definition}

\begin{thmbox}{Théorème de décohérence (H)}\mTHM
\label{thm:theorem_H}\leanTodo{PT.Stochastic.ArithmeticDecoherence}
Pour toute fonction arithmétique~$f$ non négative et toute profondeur
de crible~$K$~:
\begin{equation}
  S_{K+1}(f) \;\geq\; S_K(f).
  \label{eq:theorem_H}
\end{equation}
C'est la \textbf{seconde loi de la thermodynamique pour le crible}~:
l'entropie de la distribution des classes de sauts est non
décroissante sous le canal du crible. La convergence est
exponentielle, à un taux contrôlé par le trou spectral
$1 - |\lambda_2(T_3)| = 1/2 = \sper$ de la matrice de transfert~(T3).
\end{thmbox}

\begin{proof}
La preuve combine trois ingrédients déjà établis dans M1--M3 : (i)
la structure doublement stochastique de la matrice de transfert du
crible (T3), (ii) le théorème de majorisation de
Hardy--Littlewood--Pólya, et (iii) la concavité de Schur de
l'entropie de Shannon.

\textbf{Étape~1 : réduction à la majorisation.}
Soit $p^{(K)} = (p_c(f,K))_{c \in \Z/3\Z}$ la distribution
d'énergie de classe à la profondeur~$K$. Par construction,
l'application d'un pas supplémentaire de crible se décrit par
l'action de la matrice de transfert doublement stochastique~$T_3$
(T3, \cite{companion_M1}) sur le simplexe des distributions de
probabilité sur les $3$~classes résiduelles :
\[
  p^{(K+1)} \;=\; T_3 \cdot p^{(K)}.
\]
La double stochasticité de~$T_3$ est exacte par T3 (chaque ligne
\emph{et} chaque colonne somme à~$1$).

\textbf{Étape~2 : majorisation de Hardy--Littlewood--Pólya.}
On note $p \succ q$ la relation standard de majorisation : $q$ est
\emph{majorée par} $p$ si $\sum_{i=1}^{k} q_i^{\downarrow} \le
\sum_{i=1}^{k} p_i^{\downarrow}$ pour tout $k$, avec égalité en
$k = n$, où $\downarrow$ désigne le réordonnancement décroissant.
Le théorème de Hardy--Littlewood--Pólya~\cite{HLP1934} affirme :
pour toute matrice doublement stochastique~$M$ et toute
distribution~$p$,
\[
  p \;\succ\; M p.
\]
Appliqué à $M = T_3$ et $p = p^{(K)}$, ceci donne
$p^{(K)} \succ p^{(K+1)}$.

\textbf{Étape~3 : concavité de Schur de l'entropie de Shannon.}
L'entropie de Shannon
$S(p) = -\sum_c p_c \log p_c$
est une fonction strictement Schur-concave sur le simplexe : si
$p \succ q$, alors $S(p) \le S(q)$, avec inégalité stricte sauf si
$p$ et $q$ sont des permutations l'une de l'autre. En combinant
avec l'étape~2,
\[
  S_K(f) \;=\; S(p^{(K)}) \;\le\; S(p^{(K+1)}) \;=\; S_{K+1}(f),
\]
ce qui est~\eqref{eq:theorem_H}.

\textbf{Étape~4 : taux de convergence.}
Le taux exponentiel d'approche vers l'entropie d'équilibre
$S_{\max} = \log 3$ est gouverné par le module de la seconde plus
grande valeur propre de~$T_3$. Par T3, $\mathrm{spec}(T_3) =
\{1, -1\}$, d'où $|\lambda_2(T_3)| = 1$ pour l'antidiagonale nue.
Après normalisation en matrice stochastique par lignes
(étape~3 de la preuve de T2 dans~\cite{companion_M3}), la
seconde valeur propre non triviale devient $\lambda_2 = -1/2$, de
module $|\lambda_2| = 1/2 = \sper$. La théorie de
Perron--Frobenius fournit alors le taux exponentiel
$\gamma_{\rm dec} = 1 - |\lambda_2| = 1/2$.
\end{proof}

\begin{remarque}
Le théorème~H est inconditionnel : il n'exige aucune hypothèse
probabiliste sur~$f$, seulement la non-négativité. L'égalité
$S_{K+1}(f) = S_K(f)$ a lieu si et seulement si $p^{(K)}$ est déjà
la distribution uniforme, c'est-à-dire le point fixe de~$T_3$ ;
dans ce cas $S_K(f) = S_{\max} = \log 3$ et le crible a atteint
l'équilibre informationnel.
\end{remarque}

\paragraph{Taux de décohérence.}
Le taux $\gamma(f,K) = 1 - (S_{\max} - S_K)/(S_{\max} - S_{K-1})$
mesure la rapidité avec laquelle $f$ approche l'entropie d'équilibre
$S_{\max} = \ln 3$. Il est lié au trou spectral~:
$\gamma \sim 1 - |\lambda_2(T_3)|$.

\paragraph{Information accessible.}
La quantité $I_{\mathrm{acc}}(f,K) = S_{\max} - S_K(f) \geq 0$ est
l'information que le crible préserve encore sur~$f$ à la
profondeur~$K$. La perte d'information par pas
$\Delta I(K) = I_{\mathrm{acc}}(K) - I_{\mathrm{acc}}(K+1) \geq 0$
est non négative d'après le théorème~H.

\paragraph{Distance entropique.}
La transformée de décohérence définit une distance combinant
entropie, pureté et moyennes signées par classe~:
\begin{equation}
  d_{\mathcal{D}}(f,g)
  = \sqrt{\sum_K \bigl[
    (S_K^f - S_K^g)^2 + (P_K^f - P_K^g)^2
    + \textstyle\sum_c (m_c^f - m_c^g)^2
  \bigr]},
  \label{eq:decoherence_distance}
\end{equation}
où $P_K = \mathrm{Tr}(\rho_K^2)$ est la pureté et $m_c$ la moyenne
signée par classe. Les termes de moyenne signée garantissent que des
fonctions d'identique $|f|^2$ (par ex.\ $\mathbf{1}$ et $\chi_3$)
soient distinguées.

Script~:
\texttt{ch\_math\_structures/test\_decoherence\_transform.py}
(scripts compagnons).


\subsubsection{Comparaison des transformées PT}

\begin{center}
\small
\begin{tabular}{lccccc}
\toprule
\textbf{Propriété} & $P$ (M16) & $\mathcal{H}$ (M34) &
  $\mathcal{T}$ (M33) & $\mathcal{D}$ (M35) & Fourier \\
\midrule
Base spectrale & $v_\pm$ de $T_3$ & $\sinp{p}$ &
  $\bigotimes_q \mathrm{evecs}(T_q)$ & entropie & $e^{ikx}$ \\
Paramètre & profondeur $K$ & premiers $p$ & $(K,q_1,q_2,q_3)$ &
  profondeur $K$ & fréquence $k$ \\
Dim.\ par $K$ & 2 & $K$ (premiers) & 105 & 1 & $\infty$ \\
Linéaire & OUI & NON & OUI & NON & OUI \\
Monotone & NON & NON & NON & \textbf{OUI} & NON \\
Voit le crible & \textbf{OUI} & \textbf{OUI} & \textbf{OUI} &
  \textbf{OUI} & NON \\
Corr.\ inter-mod. & NON & partielle & \textbf{OUI} & NON & NON \\
Courbure & NON & \textbf{OUI} & NON & NON & NON \\
\bottomrule
\end{tabular}
\end{center}

Aucune de ces transformées n'est un cas particulier de Fourier ou
Mellin~: la base spectrale est celle de la matrice de transfert du
crible~$T_q$, et non~$\{e^{ikx}\}$~; le paramètre est la profondeur
du crible~$K$, et non une fréquence~; et le théorème~H n'a aucun
analogue en analyse harmonique classique.


% ============================================================
\subsection{Crible quantique et réseaux tensoriels}
\label{sec:ms_quantum}

\subsubsection{Le crible comme canal CPTP}

La matrice de transfert $T_K$ peut être interprétée comme un canal
quantique dans la représentation de Kraus. Soit $\rho_K$ la matrice
densité encodant la distribution des sauts au niveau~$K$.

\begin{thmbox}{Décohérence arithmétique}
\label{thm:decoherence}
Le canal du crible $\mathcal{E}_K\colon \rho \mapsto T_K \rho
T_K^\dagger$ est une application complètement positive préservant la
trace (CPTP) telle que~:
\begin{enumerate}[nosep]
\item L'entropie de von~Neumann $S_{\rm vN}(\rho_K)$ est strictement
  croissante avec $K$ (décohérence).
\item La pureté $\mathrm{Tr}(\rho_K^2) \to 1/3$ lorsque $K \to \infty$
  (convergence vers l'état maximalement mixte sur 3 classes).
\item Le taux de décohérence est contrôlé par $|\lambda_2(T_K)|$.
\end{enumerate}
Script~: \texttt{ch\_math\_structures/test\_quantum\_sieve.py}
(scripts compagnons).
\end{thmbox}

Ce n'est pas une affirmation selon laquelle le crible \emph{est} un
système quantique. C'est l'observation que le formalisme mathématique
de la décohérence quantique (canaux CPTP, évolution de Lindblad,
croissance de l'entropie) est le langage naturel pour décrire la
perte d'information du crible à chaque pas.

\subsubsection{Réseau tensoriel TCR}

La factorisation TCR $\mathbb{Z}/m\mathbb{Z} \cong
\prod_{p|m} \mathbb{Z}/p\mathbb{Z}$ induit une structure de produit
tensoriel sur l'espace d'états. La dynamique du crible peut s'écrire
comme un état produit de matrices (MPS)~:

\begin{thmbox}{Représentation MPS}
\label{thm:mps}
La distribution des sauts au niveau $K$ admet une représentation MPS
exacte de \textbf{dimension de liaison $\chi = 3$}. L'entropie
d'intrication entre toute bipartition des facteurs TCR satisfait
$S_{\rm ent} \leq \log 3$, et la longueur de corrélation $\xi$ est
finie ($\xi < \infty$)~: le réseau tensoriel est \textbf{gappé}.
Script~: \texttt{ch\_math\_structures/test\_tensor\_network.py}
(scripts compagnons).
\end{thmbox}

La dimension de liaison 3 est \emph{exacte}, et non une
approximation. La structure TCR du crible est efficacement capturée
par un réseau tensoriel 1D sans erreur de troncation.


% ============================================================
\subsection{Codes correcteurs d'erreurs depuis les sauts}
\label{sec:ms_codes}

\subsubsection{Prédiction multi-modulaire}

Le prédicteur multi-modulaire utilise les résidus modulo
$q = 3, 5, 7$ simultanément. L'information mutuelle entre modules
est substantielle ($\mathrm{MI} \sim 1$\,nat), confirmant que les
facteurs TCR ne sont \emph{pas} informationnellement indépendants. La
prédiction de la classe de saut suivante s'améliore de $+61{,}2\%$
par rapport au prédicteur modulo~3 seul.

Script~:
\texttt{ch\_math\_structures/test\_multi\_modular\_predictor.py}
(scripts compagnons).

\subsubsection{Code quantique TCR}

La structure TCR définit un code quantique correcteur d'erreurs avec
paramètres $[[3, 5.58, 1]]$. L'écart de Knill--Laflamme est
$\Delta_{\rm KL} = 0{,}008$ (QEC approximative, quasiment exacte).
Idée clé~: l'information mutuelle entre modules est
\emph{simultanément} la ressource pour la prédiction (direction~B) et
la capacité de correction d'erreur (direction~C). Les deux directions
sont la même structure mathématique vue sous des angles différents.

Script~: \texttt{ch\_math\_structures/test\_crt\_quantum\_code.py}
(scripts compagnons).

\textbf{Limitation honnête~:} sur les résidus entiers, la distance
du code est $d = 1$ (un code produit trivial sans contraintes
inter-modulaires). La résolution provient de l'usage des résidus de
\emph{sauts}.

\subsubsection{Le code de distance de sauts (théorème principal)}

\begin{thmbox}{Théorème de distance TCR-saut}
\label{thm:gap_distance}
Soient $q_1 < q_2 < \cdots < q_n$ les $n$ premiers premiers impairs
(les modules à la profondeur du crible $K = n+1$). On définit le
\emph{code de résidus de sauts}~$\mathcal{C}$ par~:
\[
  \mathcal{C} = \bigl\{
    (g \bmod q_1,\; g \bmod q_2,\; \ldots,\; g \bmod q_n)
    : g \in \mathrm{sauts\_distincts}(K)
  \bigr\}.
\]
Alors pour $n \geq 2$~:
\begin{enumerate}[nosep]
\item L'application TCR est injective sur les sauts~:
  $|\mathcal{C}| = |\mathrm{sauts\_distincts}|$.
\item La distance de Hamming minimale satisfait $d = n$ pour $n = 2$
  (cas frontière) et $d = n - 1$ pour $n \geq 3$.
\item Le code corrige $t = \lfloor(d-1)/2\rfloor$ erreurs.
\item Le rapport $d/n$ est croissant pour $n \geq 3$ et $d/n \to 1$
  lorsque $n \to \infty$ (asymptotiquement MDS).
\end{enumerate}
\end{thmbox}

\begin{proof}[Esquisse de preuve]
Supposons que deux sauts distincts $g_1 \neq g_2$ coïncident en $m$
positions de coordonnées $i_1, \ldots, i_m$. Alors
$q_{i_1} \cdots q_{i_m}$ divise $g_1 - g_2$. Puisque
$|g_1 - g_2| \leq \mathrm{saut}\_\mathrm{max}\sim 2p_K$ (croissance
linéaire) tandis que $q_{i_1} \cdots q_{i_m} \geq 3 \cdot 5 = 15$
pour $m \geq 2$ (croissance exponentielle), on obtient $m \leq 1$.
Donc $d \geq n - 1$. Des paires atteignant exactement $n - 1$
positions différentes existent (vérifié pour $K = 3 \ldots 7$), donc
$d = n - 1$ exactement.
\end{proof}

\begin{center}
\begin{tabular}{ccccc}
\toprule
$K$ & $n$ (modules) & $d$ (distance) & $t$ (corrigible) & $d/n$ \\
\midrule
3 & 2 & 2 & 0 & 1{,}000 \\
4 & 3 & 2 & 0 & 0{,}667 \\
5 & 4 & 3 & 1 & 0{,}750 \\
6 & 5 & 4 & 1 & 0{,}800 \\
7 & 6 & 5 & 2 & 0{,}833 \\
\bottomrule
\end{tabular}
\end{center}

\paragraph{Comparaison.}
Un code sur des résidus \emph{entiers} (et non des résidus de sauts)
a toujours $d = 1$~: c'est un code produit sans contraintes
inter-modulaires. Le code des résidus de sauts atteint $d = n-1$ car
les valeurs des sauts sont bornées (linéairement en $p_K$) tandis que
les produits des modules croissent exponentiellement. La \emph{borne
sur les sauts} est l'ingrédient qui crée la distance.

\paragraph{Premiers d'écho.}
Les premiers $p > p_K$ situés sous $\mathrm{saut}\_\mathrm{max}$
agissent comme des erreurs non détectables~: ils modifient la suite
des sauts sans changer les coordonnées du code. Lorsque $K$ augmente,
la distance du code $d = n - 1$ croît plus rapidement que le nombre
de premiers d'écho, de sorte que la capacité de correction d'erreur
du code finit par dominer.

Script~:
\texttt{ch\_math\_structures/test\_gap\_distance\_code.py}
(scripts compagnons).

\subsubsection{Synthèse~: quatre visages du crible}

Le crible est simultanément~:
\begin{enumerate}[nosep]
\item Un \textbf{algorithme} générateur de premiers (théorie des
  nombres).
\item Un \textbf{canal de décohérence} CPTP (information quantique).
\item Un \textbf{code correcteur d'erreurs} avec $d = n - 1$ (théorie
  des codes).
\item Une \textbf{contraction spectrale} de trou $1 - |\lambda_2|$
  (systèmes dynamiques).
\end{enumerate}
La distance du code $d = n - 1$ est la \emph{manifestation
arithmétique} de la persistance informationnelle~: le crible produit
un code dont la redondance croît sans borne, garantissant que le
motif des premiers est récupérable à partir d'une information
partielle.


% ============================================================
\subsection{Banc d'essai et limites honnêtes}
\label{sec:ms_benchmark}

Un banc d'essai des cinq nouvelles structures mathématiques (algèbre,
transformée, métrique, nombres, catégorie) face à des tâches
concrètes donne~:

\begin{itemize}[nosep]
\item \textbf{Classification} (premier vs.\ composé)~: $d_{\rm PT}$
  atteint $96\%$ de précision, compétitive avec les classifieurs
  fondés sur $\Omega(n)$.
\item \textbf{Prédiction}~: la prédiction multi-modulaire améliore
  modulo~3 de $+61\%$ (M30).
\item \textbf{Correction d'erreur}~: $d = n-1$ prouvé, $t \geq 1$
  pour $K \geq 5$ (M32).
\item \textbf{Détection structurelle}~: la transformée de persistance
  $P_K$ sépare proprement $\mathbf{1}$ et $\chi_3$ (M16).
\item \textbf{Cohérence catégorielle}~: les quatre foncteurs
  convergent vers la structure asymptotique (M19).
\end{itemize}

\paragraph{Limites honnêtes.}
\begin{itemize}[nosep]
\item $d_{\rm PT}$ ne remplace pas les tests de primalité (elle
  classifie statistiquement, et non déterministiquement).
\item Le code des sauts a $|\mathcal{C}| \sim 2n$ mots de code~--- trop
  peu pour la communication pratique. Sa valeur est théorique.
\item Le canal CPTP est une analogie structurelle, et non une
  réalisation physique.
\item L'extension quantique CSS donne $d_{\rm CSS} \sim 1$ pour
  $n$~grand (la distance classique ne se transfère pas pleinement).
\end{itemize}

Script~: \texttt{ch\_math\_structures/test\_capabilities.py}
(scripts compagnons).

\medskip
\noindent\textbf{Total~:} 32 scripts compagnons, tous validés
(scripts compagnons).

\section{Mécanique complexe du crible}
\label{sec:complex}

\epigraph{Le chemin le plus court entre deux vérités du domaine réel
passe par le domaine complexe.}{Jacques Hadamard}

% ---- Status Box ----------------------------------------
\begin{statusbox}
\textbf{OBJECTIF~:}
  Relever PT depuis l'observable réelle $\sinp{p}$ vers la variable
  complexe naturelle $w_p = (1 - e^{2i\theta_p})/2$, en mettant au
  jour la structure cachée~: géométrie circulaire, force holomorphe,
  moment angulaire et intégrabilité de Liouville complète.

\textbf{ENTRÉES~:}
  L'article compagnon~\cite{companion_M2} (identité d'holonomie
  $\sinp{p} = \delta_p(2-\delta_p)$), l'article
  compagnon~\cite{companion_M1} (TFS, action $S_{\mathrm{PT}}$),
  l'article compagnon~\cite{companion_physics} ($\alphaEM$ depuis le
  produit du crible).

\textbf{RÉSULTATS PRINCIPAUX~:}
  \begin{itemize}[nosep]
  \item Théorème du cercle~: tous les $w_p$ sont sur
    $\mathcal{C}(1/2,\,0;\, s)$, rayon $s = 1/2$
    (Thm~\ref{thm:circle}).
  \item $|w|^2 = \mathrm{Re}(w) = \sinp{p}$~: le module carré égale la
    projection réelle (ID).
  \item $\alpha = |W|^2 = \prod |w_p|^2$~: la structure fine comme
    module complexe (ID).
  \item $|T_{12,\mathbb{C}}| > T_{12,\mathrm{re}}$ de 29 à 62\%~:
    renforcement de la borne spectrale.
  \item Force holomorphe $F(w) = i/w - 2i$ (noyau de Cauchy, pôle en
    $w = 0$).
  \item Moment angulaire $L = \sinp{p} = |w|^2$ comme charge de
    Noether.
  \item Relation d'incertitude $|p_w|\cdot|w| = 1$ (identité exacte,
    et non une borne).
  \item Équipartition $E_p \to \kappa/2 = 1$, température géométrique
    $T = 1/s = 2$.
  \item Intégrabilité de Liouville complète sur $\mathbb{T}^n$~;
    $\alpha = \prod L_p$.
  \item Coefficient FSR QED~: $3/(4\pi) = N_{\mathrm{spatial}}\cdot
    s^2/\mathrm{Circ}(\mathcal{C})$ (DER).
  \item Facteur de boucle~: $16\pi^2 = \Omega_2^2$ à partir de
    $N_{\mathrm{spatial}} = 3$ (DER).
  \end{itemize}

\textbf{STATUT~:}
  \DERTAG\ + \textsc{[id]} (théorème du cercle, 7~identités, force
  holomorphe, intégrabilité de Liouville, corrections radiatives).
  10~scripts (M41--M50), 127/127 PASS.

\textbf{NON PROUVÉ~:}
  Interprétation physique de $\arg(W) = 0{,}937\pi$ (lien potentiel
  avec la violation~CP). Lien entre les corrections de
  Bernoulli--zêta et le catalogue NLO existant.
\end{statusbox}


% ============================================================
\subsection{La variable complexe \texorpdfstring{$w_p$}{wp}}
\label{sec:cx_variable}

À chaque premier $p$ de paramètre $q \in \{\qplus, \qminus\}$,
l'angle d'holonomie du crible $\theta_p$ définit une variable
complexe naturelle~:

\begin{equation}
  w_p \;=\; \frac{1 - e^{2i\theta_p}}{2}\,.
  \label{eq:wp_def}
\end{equation}

Les parties réelle et imaginaire se décomposent comme~:
\begin{align}
  \mathrm{Re}(w_p) &= \sin^2\theta_p = \delta_p(2-\delta_p)
    && \text{(perte = observable PT standard)}\,, \label{eq:wp_re}\\
  \mathrm{Im}(w_p) &= -\tfrac{1}{2}\sin(2\theta_p) = -\sin\theta_p\cos\theta_p
    && \text{(cohérence = observable cachée)}\,, \label{eq:wp_im}\\
  |w_p| &= \sin\theta_p\,,\quad
  |w_p|^2 = \sin^2\theta_p = \mathrm{Re}(w_p)
    && \text{(identité fondamentale)}\,. \label{eq:wp_mod}
\end{align}

PT standard n'utilise que $\mathrm{Re}(w_p)$. La partie imaginaire
$\mathrm{Im}(w_p) = -\sin\theta_p\cos\theta_p$ mesure la
\emph{cohérence} entre les canaux de perte et de conservation~--- la
moyenne géométrique des amplitudes de perte et de conservation, avec
un signe.

\medskip
\noindent\textbf{Valeurs numériques} ($q = \qplus = 13/15$, premiers
actifs)~:
\begin{center}
\begin{tabular}{cccc}
\toprule
$p$ & $\mathrm{Re}(w_p) = \sinp{p}$ & $\mathrm{Im}(w_p)$ & $|w_p|$ \\
\midrule
3 & 0{,}3340 & $-$0{,}4716 & 0{,}5780 \\
5 & 0{,}1992 & $-$0{,}3991 & 0{,}4463 \\
7 & 0{,}0528 & $-$0{,}2237 & 0{,}2298 \\
\bottomrule
\end{tabular}
\end{center}

La cohérence $|\mathrm{Im}(w_p)|$ dépasse $\mathrm{Re}(w_p)$ pour
$p = 5, 7$~: PT standard écarte la composante \emph{dominante} de la
variable complexe.

\begin{thmbox}[Théorème du cercle]
\label{thm:circle}
Pour tous les premiers $p$ et tout $q \in (0,1)$, le point $w_p$ se
trouve sur le cercle
\begin{equation}
  \mathcal{C}\colon\;
  \bigl(\mathrm{Re}(w) - \tfrac{1}{2}\bigr)^2 + \mathrm{Im}(w)^2
  = \tfrac{1}{4}\,,
  \label{eq:circle}
\end{equation}
centré en $(1/2, 0)$ et de rayon $s = 1/2$.
\end{thmbox}

\begin{proof}
Soit $x = \sin^2\theta$, $y = -\sin\theta\cos\theta$.
Alors $x^2 + y^2 = \sin^4\theta + \sin^2\theta\cos^2\theta
= \sin^2\theta(\sin^2\theta + \cos^2\theta) = \sin^2\theta = x$.
Donc $(x - 1/2)^2 + y^2 = x^2 - x + 1/4 + y^2 = 1/4$.
\end{proof}

Le cercle passe par l'origine ($\theta = 0$~: aucune perte) et par
$(1, 0)$ ($\theta = \pi/2$~: perte totale). Lorsque $K$ augmente, les
nouveaux premiers poussent $w_p$ vers l'origine~: \emph{la
contraction est une convergence radiale sur le cercle vers $w = 0$}.

\paragraph{Sept identités fondamentales.}
Les identités suivantes valent pour tous les premiers et tout
$q \in (0,1)$~:
\begin{align}
  |w_p|^2 &= \mathrm{Re}(w_p) \,,\tag{ID1}\label{id:mod_re}\\
  \mathrm{Im}(w_p)^2 &= \mathrm{Re}(w_p)\,(1 - \mathrm{Re}(w_p))\,,\tag{ID2}\label{id:im_parab}\\
  w_p &= -i\sin\theta_p\, e^{i\theta_p}\,,\tag{ID3}\label{id:polar}\\
  \arg(w_p) &= \theta_p - \pi/2\,,\tag{ID4}\label{id:arg}\\
  \bar{w}_p &= \frac{w_p}{2w_p - 1}\,,\tag{ID5}\label{id:conj}\\
  z_p &= 1 - 2w_p = e^{2i\theta_p}\,,\tag{ID6}\label{id:u1}\\
  \delta_p &= \frac{(1 - q)[p]_q}{p}\,.\tag{ID7}\label{id:qbracket}
\end{align}

L'identité~\ref{id:mod_re} est la propriété définissante du cercle
$\mathcal{C}$. L'identité~\ref{id:conj} fait du conjugué une
\emph{transformation de Möbius}, rendant la force
$F = -i/\bar{w}$ holomorphe en $w$ seul
(section~\ref{sec:cx_force}).


% ============================================================
\subsection{Observables complexes}
\label{sec:cx_observables}

\subsubsection{Produit complexe et \texorpdfstring{$\alpha$}{alpha}}

La constante de structure fine du crible provient du produit~:
\begin{equation}
  W \;=\; \prod_{p\;\mathrm{actif}} w_p\,,\qquad
  |W|^2 = \prod_{p} |w_p|^2 = \prod_{p} \sinp{p} = \alpha_{\mathrm{nu}}\,.
  \label{eq:W_product}
\end{equation}

La \emph{phase} de $W$ est une nouvelle observable, indépendante de
$|W|^2 = \alpha$~:
\begin{equation}
  \arg(W) = \sum_p \arg(w_p) = \sum_p (\theta_p - \pi/2)\,.
  \label{eq:argW}
\end{equation}
Pour $q = \qplus$, $\arg(W) \approx 0{,}937\pi$.
Ce n'est \emph{pas} une fraction simple de $\pi$~; l'approximation
rationnelle la plus proche est $15/16$.

\subsubsection{\texorpdfstring{$T_{12}$}{T12} complexe et renforcement spectral}
\label{subsec:cx_T12}

Le trou spectral $T_{12}$ admet une extension complexe via le
caractère $\chi_3$~:
\begin{equation}
  T_{12,\mathbb{C}} = \frac{1}{3}\sum_p \chi_3(p)\, w_p
  = \frac{w_7 - w_5}{3}\,,
  \label{eq:T12C}
\end{equation}
puisque $\chi_3(3) = 0$, $\chi_3(5) = -1$, $\chi_3(7) = +1$. Le
module satisfait~:
\[
  |T_{12,\mathbb{C}}| > |T_{12,\mathrm{re}}|\,,\qquad
  \text{rapport } \frac{|T_{12,\mathbb{C}}|}{|T_{12,\mathrm{re}}|}
  = \frac{1}{\sin\theta_{\mathrm{eff}}} \in [1{,}29~;\; 1{,}62]\,.
\]
Pour $\qplus$~: $|T_{12,\mathbb{C}}| = 0{,}00920$ contre
$T_{12,\mathrm{re}} = 0{,}00712$ ($+29\%$).
Pour $\qminus$~: gain $+62\%$.
Cela renforce les bornes spectrales de la preuve de convergence T5
(article compagnon~\cite{companion_M3}) sans changer la structure de
la preuve.

\subsubsection{Action complexe}

L'action effective PT $S_{\mathrm{PT}} = -\sum_p \ln\sinp{p}$ s'étend
en~:
\begin{equation}
  S_{\mathbb{C}} = -\sum_p \ln w_p
  = -\sum_p \bigl[\ln|w_p| + i\arg(w_p)\bigr]\,,
  \label{eq:S_complex}
\end{equation}
avec $\mathrm{Re}(S_{\mathbb{C}}) = S_{\mathrm{PT}}/2$ (puisque
$\ln\sinp{p} = 2\ln|w_p|$). La partie imaginaire encode la phase
accumulée le long du chemin des premiers.


% ============================================================
\subsection{Structure algébrique profonde}
\label{sec:cx_algebra}

\subsubsection{Correspondance U(1)}

L'application affine $z_p = 1 - 2w_p = e^{2i\theta_p}$ envoie le
cercle $\mathcal{C}$ sur le cercle unité $\mathbb{S}^1$. Cela
identifie l'espace d'états du crible avec $\mathrm{U}(1)$.

\subsubsection{Fisher = Fubini--Study}

La métrique de Fisher sur l'espace de paramètres de Bernoulli se
retire en quatre fois la métrique de longueur d'arc sur
$\mathcal{C}$~:
\begin{equation}
  \mathrm{d}s^2_{\mathrm{Fisher}}
  = \frac{4}{\sinp{p}(1 - \sinp{p})}\,(\mathrm{d}\sinp{p})^2
  = 4\,\mathrm{d}\theta_p^2\,.
  \label{eq:fisher_fubini}
\end{equation}
C'est la métrique de Fubini--Study sur $\mathbb{CP}^1$. Le facteur~4
égale l'inverse du carré du rayon du cercle~: $4 = 1/s^2$.

\subsubsection{Correspondance avec la matrice densité}
\label{subsec:cx_density}

Chaque $w_p$ encode la première colonne d'une matrice densité de
rang~1 $\rho_p = |\psi_p\rangle\langle\psi_p|$ pour un système à
2~niveaux $|\psi_p\rangle = \sin\theta_p\,|0\rangle +
\cos\theta_p\,|1\rangle$~:
\begin{equation}
  w_p = (\rho_p)_{00} - i\,(\rho_p)_{01}
  = \sin^2\theta_p - i\sin\theta_p\cos\theta_p\,.
  \label{eq:density}
\end{equation}

La partie réelle est l'élément diagonal (probabilité), la partie
imaginaire l'élément hors diagonale (cohérence). Le cercle
$\mathcal{C}$ est le lieu des états purs sur le méridien
$\Phi = 0$ de la sphère de Bloch.

\subsubsection{Birapports}

Pour quatre points co-cycliques sur $\mathcal{C}$, le birapport est
\emph{exactement réel} (partie imaginaire $< 10^{-16}$). Pour les
premiers actifs~:
\[
  \mathrm{CR}(3, 5;\, 7, \infty) \approx 2{,}000
  \quad (\text{erreur } 1{,}5 \times 10^{-4})\,.
\]
Les birapports quasi-entiers reflètent la structure harmonique des
trois premiers actifs.


% ============================================================
\subsection{Corrections complexes}
\label{sec:cx_corrections}

\subsubsection{Le rayon du cercle encode \texorpdfstring{$s$}{s}}

Le cercle $\mathcal{C}$ a pour rayon $s = 1/2$. Le paramètre
fondamental de PT est donc encodé dans la géométrie du plan
complexe~: \emph{la géométrie est la théorie}.

\subsubsection{Décomposition structurelle de \texorpdfstring{$1/\alpha$}{1/alpha}}
\label{subsec:cx_bernoulli}

Le produit nu $1/\alpha_{\mathrm{nu}} = \prod(1/\sinp{p})$
se décompose via $\sin^2 = \theta^2\cdot(\sin\theta/\theta)^2$~:
\begin{equation}
  \frac{1}{\alpha_{\mathrm{nu}}}
  = \underbrace{\prod_p \frac{1}{\theta_p^2}}_{= 110{,}34\;\text{(classique)}}
  \;\times\;
  \underbrace{\prod_p \Bigl(\frac{\theta_p}{\sin\theta_p}\Bigr)^2}_{= 1{,}235\;\text{(Bernoulli--zêta)}}\,.
  \label{eq:alpha_decomp}
\end{equation}

Le facteur «~classique~» n'utilise que l'ordre dominant
$\sin\theta \approx \theta$. Le facteur Bernoulli--zêta encode les
corrections d'ordre supérieur via
$\theta/\sin\theta = 1 + \theta^2/6 + 7\theta^4/360 + \cdots$,
où les coefficients font intervenir les nombres de Bernoulli
$B_{2k}$ et implicitement $\zeta(2k)$.

\emph{Cette décomposition est une identité}~--- le facteur de
Bernoulli est déjà contenu dans $\sin^2$ et ne constitue pas une
nouvelle correction à $\alpha$. Elle révèle la structure interne du
produit nu.

\subsubsection{NLO à partir de la cohérence}

Le rapport des parties imaginaire et réelle donne le rapport
NLO/LO~:
\begin{equation}
  \frac{|\mathrm{Im}(w_p)|}{\mathrm{Re}(w_p)}
  = |\cot\theta_p|
  \sim \sqrt{p/2}\,.
  \label{eq:nlo_cot}
\end{equation}

Pour tous les premiers au-delà de $p = 3$, la composante de cohérence
(imaginaire) dépasse la composante de perte (réelle)~: le NLO
\emph{domine} le LO dans le plan complexe.


\subsubsection{Corrections radiatives depuis la géométrie circulaire}
\label{subsec:cx_radiative}

Les deux dernières formules de TQC utilisées dans le catalogue NLO
(article compagnon~\cite{companion_physics}, annexe~B)~--- le
coefficient QED de rayonnement d'état final et le facteur de vertex
à une boucle~--- se décomposent entièrement en quantités PT-natives.

\begin{derbox}{FSR QED depuis la géométrie circulaire}
\label{der:fsr_circle}
La correction QED de rayonnement d'état final pour un fermion de
charge~$Q_f$ est~:
\begin{equation}
  \delta_{\mathrm{QED}} = \frac{3\alpha Q_f^2}{4\pi}
  = \frac{N_{\mathrm{spatial}}\cdot s^2}{\mathrm{Circ}(\mathcal{C})}
  \;\cdot\;\alpha\,Q_f^2\,,
  \label{eq:fsr_pt}
\end{equation}
où $N_{\mathrm{spatial}} = 3$ (théorème, R38b),
$s = 1/2$ (fondamental), et
$\mathrm{Circ}(\mathcal{C}) = 2\pi s = \pi$.
Le coefficient $3/4 = N_{\mathrm{spatial}}\cdot s^2$ est donc
déterminé par le nombre de dimensions spatiales et le rayon du crible.
\end{derbox}

Interprétation physique~: $N_{\mathrm{spatial}}$ compte les directions
de polarisation indépendantes pour le quantum émis~;
$s^2 = 1/4$ est le rapport $\mathrm{Aire}(\mathcal{C})/\pi$~;
et $\mathrm{Circ}(\mathcal{C}) = \pi$ normalise l'espace des phases
d'émission sur le cercle.

\begin{derbox}{Facteur de boucle depuis l'angle solide}
\label{der:loop_factor}
Le facteur de normalisation à une boucle est~:
\begin{equation}
  16\pi^2 = \Omega_2^2 = \biggl(
  \frac{2\pi^{N_{\mathrm{spatial}}/2}}{\Gamma(N_{\mathrm{spatial}}/2)}
  \biggr)^{\!2},
  \label{eq:loop_factor}
\end{equation}
où $\Omega_2 = 4\pi$ est l'angle solide de la sphère unité $S^2$ en
$N_{\mathrm{spatial}} = 3$ dimensions spatiales.
\end{derbox}

La correction de vertex
$\rho_b = 1 - G_F m_t^2/(4\sqrt{2}\,\pi^2)$
(R43, non universelle $Z\to b\bar{b}$) se réécrit sous forme de
Yukawa~:
\begin{equation}
  \rho_b = 1 - \frac{y_t^2}{\Omega_2^2}\,,
  \qquad y_t = \frac{\sqrt{2}\,m_t}{v}\,,
  \label{eq:rho_b_pt}
\end{equation}
où chaque quantité est PT-dérivée~: $y_t$ depuis la hiérarchie des
masses de quarks, $v$ (R51), et $\Omega_2^2 = 16\pi^2$ depuis
$N_{\mathrm{spatial}} = 3$. La valeur numérique $\rho_b = 0{,}9938$
correspond exactement au résultat standard (identité algébrique).

\paragraph{Intégrale de contour.}
La force holomorphe $F(w) = i/w - 2i$ satisfait~:
\begin{equation}
  \oint_{\mathcal{C}} F\,\mathrm{d}w = -\pi = -\mathrm{Circ}(\mathcal{C})\,.
  \label{eq:contour_F}
\end{equation}
Puisque $w = 0$ se trouve \emph{sur} $\mathcal{C}$ (pôle au bord), la
formule de Sokhotski--Plemelj donne
$\oint(i/w)\,\mathrm{d}w = \pi i \cdot i = -\pi$ (la moitié du résidu
plein). Le travail total effectué par la force le long d'un parcours
de $\mathcal{C}$ égale sa circonférence~: une identité topologique
reliant le rayonnement à la géométrie.

\paragraph{Statut.}
Avec ces deux dérivations,
\emph{toutes les corrections NLO/NNLO utilisées dans l'article
compagnon}~\cite{companion_physics} \emph{sont désormais
PT-dérivées}~;
plus aucune formule TQC externe ne subsiste dans la chaîne
d'habillage. (Script~: M50,
\texttt{test\_complex\_radiative}, 18/18~PASS.)


% ============================================================
\subsection{Mécanique holomorphe}
\label{sec:cx_force}

\subsubsection{La force}

On définit le moment complexe par
$p_w = \mathrm{d}S_{\mathbb{C}}/\mathrm{d}\theta = -\cot\theta - i$.
La force associée dans l'espace $w$ est~:
\begin{equation}
  F(w) = \frac{i}{w} - 2i\,.
  \label{eq:force}
\end{equation}

C'est une fonction holomorphe sur $\mathbb{C} \setminus \{0\}$ avec
un pôle simple en $w = 0$ et un résidu $\mathrm{Res}(F, 0) = i$.
L'identité~\ref{id:conj} $\bar{w} = w/(2w-1)$ convertit
$F = -i/\bar{w}$ en la forme holomorphe ci-dessus.

\paragraph{Identité de rotation.}
\begin{equation}
  F(w_p)\cdot\mathrm{Re}(w_p) = -i\,w_p\,.
  \label{eq:rotation}
\end{equation}
La force multipliée par l'observable égale la variable complexe
tournée de $-\pi/2$. La multiplication par $\sinp{p}$ effectue une
rotation d'un quart de tour.

\paragraph{Force azimutale universelle.}
\begin{equation}
  \mathrm{Im}\!\Bigl(\frac{\mathrm{d}S_{\mathbb{C}}}{\mathrm{d}\theta}\Bigr)
  = -1 = -\frac{1}{2s}\,,
  \label{eq:universal}
\end{equation}
constante pour tous les premiers. La force azimutale égale l'inverse
du diamètre du cercle $2s = 1$.

\subsubsection{Moment angulaire comme charge de Noether}

Le moment angulaire associé à la rotation $\mathrm{U}(1)$
$w \mapsto e^{i\alpha}w$ est~:
\begin{equation}
  L_p = \mathrm{Im}(\bar{w}_p\,\mathrm{d}w_p/\mathrm{d}\theta)
  = \sinp{p} = |w_p|^2\,.
  \label{eq:angular_momentum}
\end{equation}

\emph{L'observable PT EST le moment angulaire.}

\subsubsection{Relation d'incertitude}
\label{subsec:cx_uncertainty}

\begin{equation}
  |p_w| \cdot |w_p| = 1 \qquad\text{pour tout } p\,.
  \label{eq:uncertainty}
\end{equation}

C'est une \emph{identité exacte} (et non une borne), de constante
égale au diamètre $2s = 1$ du cercle $\mathcal{C}$.
Depuis~\eqref{eq:uncertainty}~:
$\sinp{p} = |w|^2 = 1/(1 + \cot^2\theta)
= 1/(1 + \mathrm{Re}(p_w)^2)$,
retrouvant la règle de Born depuis le moment.


% ============================================================
\subsection{Intégrabilité et variables action--angle}
\label{sec:cx_integrability}

\subsubsection{Actions de Liouville}

On définit la variable d'action et le hamiltonien~:
\begin{equation}
  L_p = \sinp{p}\,,\qquad
  H_p = -\ln\sin\theta_p = -\tfrac{1}{2}\ln L_p\,.
  \label{eq:liouville}
\end{equation}

Le crochet de Poisson $\{L_p, H_{p'}\} = 0$ pour $p \neq p'$
(indépendance des premiers), et $\{L_p, H_p\} = 0$ car les deux ne
dépendent que de $\theta_p$ (un degré de liberté par premier). Le
système est \emph{complètement intégrable} au sens de Liouville, et
vit sur un tore $\mathbb{T}^n$.

La constante de structure fine est le produit des actions de
Liouville~:
\begin{equation}
  \alpha = \prod_{p\;\mathrm{actif}} L_p
  = \prod_{p\;\mathrm{actif}} \sinp{p}\,.
  \label{eq:alpha_liouville}
\end{equation}

\subsubsection{Équipartition}
\label{subsec:cx_equipartition}

On définit l'énergie cinétique $E_p = \frac{1}{2}p\,\theta_p^2$.
Pour $p$ grand, $\theta_p \sim \sqrt{2/p}$, donc~:
\begin{equation}
  E_p \;\to\; \frac{\kappa}{2} = 1\,,\qquad
  T_{\mathrm{geom}} = \kappa = \frac{1}{s} = 2\,.
  \label{eq:equipartition}
\end{equation}

Correction exacte~: $E_p = 1 - 2/p + O(1/p^2)$.
La température géométrique $T = 2$ égale la courbure du cercle
$\mathcal{C}$ (inverse du rayon $s = 1/2$).

\subsubsection{Corrections de Bernoulli--zêta}

Le rapport $\theta/\sin\theta$ admet le développement de Bernoulli~:
\begin{equation}
  \frac{\theta}{\sin\theta}
  = 1 + \frac{\zeta(2)}{\pi^2}\,\theta^2
    + \frac{7\zeta(4)}{2\pi^4}\,\theta^4
    + \cdots
  = \sum_{k=0}^{\infty} \frac{(-1)^{k+1}(2^{2k}-2)\,B_{2k}}{(2k)!}\,\theta^{2k}\,.
  \label{eq:bernoulli}
\end{equation}

Ceci relie le facteur de «~correction quantique~»
$\prod(\theta/\sin\theta)^2 = 1{,}235$ dans~\eqref{eq:alpha_decomp}
aux valeurs de la fonction zêta de Riemann aux entiers pairs.

\subsubsection{Régularisation UV}

La discrétisation des premiers fournit une coupure UV naturelle~:
$p_{\min} = 3$ (le plus petit premier actif). Les premiers d'écho
($p = 11, 13$) servent de contre-termes dans l'action effective,
analogues aux régulateurs de Pauli--Villars. Il n'y a \emph{ni
divergences, ni pôles, ni dépendance à la coupure}~: la mécanique
complexe est finie par construction.


% ============================================================
\subsection{Synthèse}
\label{sec:cx_synthesis}

\subsubsection{PT réelle vs.\ PT complexe}

\begin{center}
\begin{tabular}{lll}
\toprule
\textbf{Concept} & \textbf{PT réelle} & \textbf{PT complexe} \\
\midrule
Observable      & $\sinp{p} \in [0,1]$    & $w_p \in \mathcal{C} \subset \mathbb{C}$ \\
Géométrie       & Intervalle $[0,1]$      & Cercle de rayon $s = 1/2$ \\
Métrique        & Fisher $4\,\mathrm{d}\theta^2$ & Fubini--Study sur $\mathbb{CP}^1$ \\
Cohérence       & Écartée                 & $\mathrm{Im}(w_p) = -\sin\theta\cos\theta$ \\
$\alpha$        & $\prod\sinp{p}$         & $|W|^2$, plus la phase $\arg(W)$ \\
Action          & $S = -\sum\ln\sinp{p}$  & $S_{\mathbb{C}} = -\sum\ln w_p$ \\
Force           & $-\mathrm{d}S/\mathrm{d}\theta$ (réelle) & $F(w) = i/w - 2i$ (holomorphe) \\
Mom.\ angulaire & Non défini              & $L = \sinp{p}$ (charge de Noether) \\
Incertitude     & Non définie             & $|p|\cdot|w| = 1$ (identité) \\
Intégrabilité   & Non manifeste           & Liouville sur $\mathbb{T}^n$ \\
Température     & $\beta = 1/\mu$         & $T = 1/s = 2$ (courbure) \\
\bottomrule
\end{tabular}
\end{center}

\subsubsection{Impact sur les preuves existantes}

\begin{enumerate}[nosep]
\item \textbf{Bornes spectrales} (article
  compagnon~\cite{companion_M3})~:
  $|T_{12,\mathbb{C}}| > T_{12,\mathrm{re}}$ resserre le trou
  spectral de 29 à 62\%. Les preuves T4/T5 restent valides et sont
  renforcées.
\item \textbf{Structure fine} (article
  compagnon~\cite{companion_physics})~:
  $\alpha = |W|^2$ fournit une interprétation géométrique~; la
  décomposition structurelle~\eqref{eq:alpha_decomp} révèle le
  contenu Bernoulli--zêta déjà présent dans $\sinp{p}$. Aucun
  changement numérique.
\item \textbf{Métrique de Fisher} (article
  compagnon~\cite{companion_M2})~:
  La métrique $4\,\mathrm{d}\theta^2$ est Fubini--Study sur le
  cercle. C'est une reformulation, et non une correction.
\item \textbf{Règle de Born} (article
  compagnon~\cite{companion_physics})~:
  $\sinp{p} = |w|^2$ est littéralement $|\psi|^2$, et la
  correspondance avec la matrice densité~\eqref{eq:density} fournit
  un plongement quantique explicite.
\end{enumerate}

\subsubsection{Directions ouvertes}
\label{subsec:cx_open}

\begin{enumerate}[nosep]
\item Signification physique de $\arg(W) = 0{,}937\pi$~: lien avec
  l'invariant de Jarlskog ou la violation~CP~?
\item Accumulation de phase de Berry lorsque $K$ augmente~:
  quantification de
  $\gamma_K = \sum \mathrm{Im}\langle\psi_k|\psi_{k+1}\rangle$~?
\item Autres processus radiatifs dérivables depuis l'intégrale de
  contour $\oint_{\mathcal{C}} F\,\mathrm{d}w = -\pi$ et la structure
  de noyau de Cauchy de~$F$.
\item $q$-déformation et brisure du cercle~: lorsque $q \to 1$, le
  cercle s'effondre en un point~; lorsque $q \to 0$, il s'étend au
  disque complet.
\end{enumerate}


% ============================================================
% Script references
% ============================================================
\paragraph{Scripts.}
M41~(\texttt{test\_imaginary\_PT}),
M42~(\texttt{test\_complex\_plane}),
M43~(\texttt{test\_complex\_deep}),
M44~(\texttt{test\_complex\_directions}),
M45~(\texttt{test\_complex\_corrections}),
M46~(\texttt{test\_complex\_questions}),
M47~(\texttt{test\_complex\_force}),
M48~(\texttt{test\_complex\_mechanics}),
M49~(\texttt{test\_complex\_final}),
M50~(\texttt{test\_complex\_radiative}).
Tous dans \texttt{scripts/ch\_math\_structures/}.


\newpage

\section{Mathématiques prédictives}
\label{sec:PM}

\epigraph{Donnez-moi un champ et une cascade de contraintes. Je vous
dirai combien de degrés de liberté survivent à chaque niveau, quel
est leur coût d'activation, et quelle géométrie ils portent.}{Le
programme MP}

% ---- Status Box (R1) ----------------------------------------
\begin{statusbox}
\textbf{OBJECTIF~:}
  Développer le cadre des \emph{mathématiques prédictives}~(MP)~: un
  langage universel pour prédire comment les cascades de contraintes
  façonnent les degrés de liberté qui survivent. Instancier le cadre
  sur le crible d'Ératosthène et dériver des théorèmes en forme close
  pour la dimension de couche, le spectre d'activation et le seuil
  d'activation, y compris le mécanisme de rang fantôme.

\textbf{ENTRÉES~:}
  Angles d'holonomie $\sinp{p}$ et dimensions anormales $\gamp{p}$
  (article compagnon~\cite{companion_M2}), point fixe $\mustar = 15$
  (article compagnon~\cite{companion_M3}), $N_{\mathrm{spatial}} = 3$
  (article compagnon~\cite{companion_M3}), factorisation TCR (article
  compagnon~\cite{companion_M1}).

\textbf{RÉSULTATS PRINCIPAUX~:}
  \begin{itemize}[nosep]
  \item L1~: chaque premier ajoute exactement 1~DDL
    (théorème~\ref{thm:L1}, inconditionnel \THMTAG).
  \item $\dim(d) = P(d{+}1) - (2d{+}1)$, formule de dimension de
    couche (théorème~\ref{thm:dim_formula}, \THMTAG\ sous position
    générale).
  \item $\Sigma_k = \binom{4}{k}$ pour $k \ge \mathrm{AT}$, spectre
    d'activation (théorème~\ref{thm:sigma}, \THMTAG\ sous position
    générale).
  \item Formule AT en trois régimes~: écho, construction,
    asymptotique (théorèmes~\ref{thm:AT_A}--\ref{thm:AT_C},
    \THMTAG).
  \item Identité de partition~:
    $\mathrm{fantome}(1) + \mathrm{struct}(1) = n_{\mathrm{base}}$
    (identité~\ref{id:C1}, \IDTAG).
  \item Conjecture MP--PT prouvée~: $\mathrm{AT} = 1$ pour tout
    $p \ge 41$ (corollaire~\ref{cor:PMPT}).
  \end{itemize}

\textbf{STATUT~:}
  \THMTAG~(L1, AT parties A--B) + \THMTAG~(dim, $\Sigma$, AT
  partie~C, sous position générale) + \IDTAG~(C1)

\textbf{NON PROUVÉ~:}
  La conjecture de position générale
  (conjecture~\ref{hyp:gen_pos}) est vérifiée numériquement pour les
  13~strates (11--59) mais n'est pas dérivée à partir de premiers
  principes arithmétiques. La seconde instanciation (codes
  correcteurs d'erreurs) est \VALTAG, et non \THMTAG.
\end{statusbox}

% ============================================================
\subsection{Le cadre MP}
\label{sec:PM_framework}

\subsubsection{Motivation}
\mTHM

La théorie de la persistance, telle que développée dans les articles
compagnons~\cite{companion_M1,companion_M2,companion_M3}, prouve que
le crible d'Ératosthène converge vers un unique point fixe
$\mustar = 15$ avec premiers actifs $\{3, 5, 7\}$. Les axiomes de
pont (article compagnon~\cite{companion_M5}) envoient ensuite cette
structure sur la physique. Mais une question reste intacte~:

\begin{quote}
\emph{Que se passe-t-il, structurellement, lorsqu'un premier est
ajouté à la base du crible~?}
\end{quote}

Le cadre des \emph{mathématiques prédictives}~(MP) répond à cette
question. Ce n'est pas une mathématique des objets~--- il prédit ce
qui \emph{devient inévitable} lorsque des contraintes sont empilées
sur un champ brut. Trois sorties~:
\begin{enumerate}[nosep]
  \item le nombre de degrés de liberté (DDL) survivants à chaque
    niveau de contrainte,
  \item le coût d'activation pour débloquer le DDL suivant,
  \item la géométrie portée par les survivants (spectre de
    transversalité, métrique de Fisher, invariants de persistance).
\end{enumerate}

\subsubsection{Définitions universelles}
\label{sec:PM_defs}

\begin{definition}[Cascade de contraintes]
\label{def:cascade}
Une \emph{cascade de contraintes} est un tuple $(F, C, S, G, I)$ où
$F$ est un champ brut de candidats,
$C = (C_1, C_2, \ldots)$ est une famille ordonnée de contraintes,
$S_d = \{x \in F : C_1(x) \wedge \cdots \wedge C_d(x)\}$ est
l'ensemble des survivants à la profondeur~$d$,
$G$ est la géométrie lisible depuis la distribution des survivants,
et $I$ est une famille d'invariants stables sous la cascade.
\end{definition}

\begin{definition}[Matrice d'observation]
\label{def:obs_matrix}
Étant donné une famille de sondes $\eta = (\eta_1, \ldots, \eta_n)$ et
une profondeur~$d$, la \emph{matrice d'observation} $\mathbf{O}(\eta,
d)$ est la matrice de Gram centrée dont les lignes sont les
espérances conditionnelles
$\mathbb{E}[\eta_j \mid \text{niveau de conditionnement } d]$
pour chaque résidu de conditionnement. La \emph{matrice historique}
à la profondeur~$d$ est la pile verticale
\[
  \mathbf{H}(\eta, d) = \begin{pmatrix}
    \mathbf{O}(\eta, 1) \\ \vdots \\ \mathbf{O}(\eta, d)
  \end{pmatrix}.
\]
Le \emph{rang cumulatif} est
$R(\eta, d) = \operatorname{rang} \mathbf{H}(\eta, d)$.
\end{definition}

\begin{definition}[Dimension de couche et seuil d'activation]
\label{def:layer_AT}
La \emph{dimension de couche} à la profondeur~$d$ est
$\dim(d) = R(\eta, d) - R(\eta, d{-}1)$.
La \emph{profondeur future} d'une strate~$p$ est le plus grand~$d$
tel que $\dim(d) > 0$.
Le \emph{seuil d'activation} $\mathrm{AT}$ est le nombre minimal de
sondes-cibles requises pour détecter le nouveau DDL à la profondeur
future.
\end{definition}

\begin{definition}[Spectre de transversalité]
\label{def:sigma}
Pour un paquet de sondes de taille~$n_p$, le \emph{spectre de
transversalité} $(\Sigma_1, \ldots, \Sigma_{n_p})$ compte combien de
sous-ensembles de taille~$k$ détectent le nouveau DDL~:
$\Sigma_k = \#\{S \subseteq \text{paquet} : |S| = k,\;
S \text{ détecte}\}$.
\end{definition}


% ============================================================
\subsection{Instanciation crible}
\label{sec:PM_sieve}

\subsubsection{Cadre}

Dans l'instanciation crible~:
\begin{itemize}[nosep]
  \item $F = \{n \in \mathbb{N} : n \ge 2\}$ (entiers candidats).
  \item $C_k$~: divisibilité par le $k$-ième premier $p_k$ (crible
    d'Ératosthène).
  \item $S_d$~: les entiers $d$-rugueux (premiers entre eux avec
    $p_1, \ldots, p_d$).
  \item Sondes~: indicateurs centrés
    $\eta(p, t)(g) = \delta(g \equiv t \bmod p) - 1/p$
    pour chaque premier~$p$ et résidu $t \in \{p{-}4, \ldots, p{-}1\}$
    (un paquet de 4~sondes par strate).
  \item Conditionnement~: résidus modulo les premiers actifs
    $\{3, 5, 7\}$ et les premiers d'écho $\{11, 13, \ldots\}$ déjà
    dans la base.
\end{itemize}

Chaque premier $p > 7$ ajouté à la base du crible définit une
\emph{strate}. La strate explore comment les invariants MP~---
profondeur future, seuil d'activation, spectre de transversalité~---
changent lorsque la contrainte~$C_p$ est imposée.

\subsubsection{Horloge primorielle}
\label{sec:PM_clock}

Les strates pavent le cycle primoriel $\Z/30\Z$, où $30 = 2\mustar$
est le primoriel des premiers actifs. Les 8~résidus copremiers
$(\Z/30\Z)^* = \{1, 7, 11, 13, 17, 19, 23, 29\}$ sont visités dans
l'ordre par les premiers $> 7$. Le premier cycle complet (strates
11--37) couvre les 8~résidus~; le cycle~2 commence à la strate~41.

\subsubsection{Données~: 13 strates explorées}
\label{sec:PM_data}

\begin{table}[htbp]
\centering
\caption{Invariants MP pour les 13 strates explorées.}
\label{tab:PM_shells}
\begin{tabular}{rrrrrl}
\toprule
Strate & Profondeur & Incr. & AT & Spectre $(\Sigma_1,\ldots,\Sigma_4)$ & Cycle \\
\midrule
11 &  4 & 1 & 1 & $(4,\;6,\;4,\;1)$ & 1 \\
13 &  5 & 1 & 1 & $(4,\;6,\;4,\;1)$ & 1 \\
17 &  6 & 1 & 2 & $(0,\;6,\;4,\;1)$ & 1 \\
19 &  7 & 1 & 2 & $(0,\;6,\;4,\;1)$ & 1 \\
23 &  8 & 1 & 2 & $(0,\;6,\;4,\;1)$ & 1 \\
29 &  9 & 1 & 3 & $(0,\;0,\;4,\;1)$ & 1 \\
31 & 10 & 1 & 3 & $(0,\;0,\;4,\;1)$ & 1 \\
37 & 11 & 1 & 3 & $(0,\;0,\;4,\;1)$ & 1 \\
41 & 12 & 1 & 1 & $(4,\;6,\;4,\;1)$ & 2 \\
43 & 13 & 1 & 1 & $(4,\;6,\;4,\;1)$ & 2 \\
47 & 14 & 1 & 1 & $(4,\;6,\;4,\;1)$ & 2 \\
53 & 15 & 1 & 1 & $(4,\;6,\;4,\;1)$ & 2 \\
59 & 16 & 1 & 1 & $(4,\;6,\;4,\;1)$ & 2 \\
\bottomrule
\end{tabular}
\end{table}

Trois faits sont visibles~:
\begin{enumerate}[nosep]
  \item L'incrément vaut toujours~1 (L1).
  \item Le spectre vaut toujours $\binom{4}{k}$ pour
    $k \ge \mathrm{AT}$, zéro en dessous.
  \item Le AT suit un motif d'\emph{escalator} au cycle~1
    (AT = 1, 1, 2, 2, 2, 3, 3, 3) et se réinitialise à AT~=~1 pour
    tout le cycle~2.
\end{enumerate}


% ============================================================
\subsection{Théorème~L1~: incrément unité}
\label{sec:PM_L1}

\begin{thmbox}{L1 --- Théorème d'incrément unité}
\label{thm:L1}
Pour tout premier $p > 7$, l'ajout de $p$ à la base du crible
augmente la profondeur future d'exactement~$1$~:
\[
  \mathrm{PF}(\text{strate } p) = \mathrm{prof\_base} + 1.
\]
\end{thmbox}

La preuve repose sur deux lemmes algébriques. \mTHM

\begin{lemme}[Saturation d'indicateurs]
\label{lem:indicator_sat}
Pour tout premier~$p$ et tous résidus distincts
$t_1 \ne t_2 \in \Z/p\Z$,
\[
  \delta(g \equiv t_1) \cdot \delta(g \equiv t_2) = 0
  \qquad\text{(identiquement, et non statistiquement).}
\]
De plus $\delta(g \equiv t)^2 = \delta(g \equiv t)$ (idempotence).
\end{lemme}

\begin{proof}
Un saut $g$ est congru à exactement un résidu modulo~$p$, donc le
produit d'indicateurs pour des résidus distincts s'annule
identiquement. L'idempotence découle de $0^2 = 0$, $1^2 = 1$.
\end{proof}

\noindent
\textbf{Conséquence.} Tout monôme de degré $\ge 2$ en les sondes
d'un \emph{seul} premier est une combinaison linéaire des sondes de
degré~1 (plus des constantes). L'algèbre des sondes pour le
premier~$p$ est engendrée par les $p{-}1$ indicateurs centrés
$\{\eta(p, t)\}_{t \ne 0}$.

\begin{lemme}[Irréductibilité TCR]
\label{lem:CRT_irred}
Pour des premiers distincts $p \ne q$, le produit
$\delta(g \equiv t \bmod p) \cdot \delta(g \equiv s \bmod q)$ égale
$\delta(g \equiv r \bmod pq)$ où~$r$ est l'unique relèvement~TCR.
Cette observable au module~$pq$ n'est \emph{pas} dans
l'engendrement des observables aux modules~$p$ et~$q$ séparément.
\end{lemme}

\begin{proof}
Supposons que
$\delta(g \equiv r \bmod pq) = \sum_i a_i f_i(g \bmod p) + \sum_j b_j
h_j(g \bmod q) + c$.
L'évaluation en $g = r$ et $g = r + p$ (même résidu mod~$p$,
différents mod~$q$) impose $b_j = 0$ pour tout~$j$. Par symétrie,
$a_i = 0$ pour tout~$i$. Mais $\delta(g \equiv r \bmod pq)$ est non
constant~: contradiction.
\end{proof}

\begin{proof}[Preuve du théorème~\ref{thm:L1}]
\textbf{Partie~I~: $\mathrm{PF} \ge \mathrm{prof\_base} + 1$.}\quad
À la profondeur $D = d{+}1$ (après ajout de $p = p_{d+1}$), les
observables incluent les produits croisés
$\eta(p_1, t_1) \cdots \eta(p_d, t_d) \cdot \eta(p, t)$~--- un
produit de $d{+}1$ sondes pour des premiers \emph{distincts}. Par
TCR itéré ($\Z/m'\Z \cong \prod \Z/p_i\Z$), ce produit contient
l'observable irréductible $\delta(g \equiv r \bmod m')$ au module
$m' = p_1 \cdots p_d \cdot p$. Puisque $\DKL(P_p \| U_p) > 0$ (le
biais des sauts mod~$p$ est non trivial d'après~T0), la composante
$\Z/p\Z$ porte de l'information non nulle. Le rang à la profondeur
$d{+}1$ est strictement plus grand qu'à la profondeur~$d$.

\textbf{Partie~II~: $\mathrm{PF} \le \mathrm{prof\_base} + 1$.}\quad
À la profondeur $D = d{+}2$, tout monôme de degré $d{+}2$ en les
sondes de $d{+}1$ premiers distincts doit faire intervenir au moins
un premier au carré (par tiroirs). Par le
lemme~\ref{lem:indicator_sat}, le facteur carré se réduit au
degré~$\le 1$, ce qui ramène le monôme au degré $\le d{+}1$. Aucune
nouvelle direction n'apparaît à la profondeur $d{+}2$.
\end{proof}

\begin{remarque}
L1 est inconditionnel~: aucune hypothèse probabiliste, aucune
hypothèse de position générale. Il vaut pour \emph{toute} suite de
sauts, pas seulement les sauts de premiers. Le théorème est une
propriété de l'algèbre TCR des indicateurs modulaires.
\end{remarque}


% ============================================================
\subsection{Formule de dimension de couche}
\label{sec:PM_dim}

\begin{conjecture}[Position générale]
\label{hyp:gen_pos}
Les corrections d'interaction entre premiers distincts dans la
matrice d'observation sont en position générale~: aucune relation
linéaire accidentelle ne survient entre les termes croisés TCR.
(Vérifié pour 13~strates, non prouvé.)
\end{conjecture}

\begin{thmbox}{Formule de dimension}
\label{thm:dim_formula}
Sous la conjecture~\ref{hyp:gen_pos}, la dimension stabilisée de la
couche à la profondeur~$d$ est
\[
  \dim(d) = P(d{+}1) - (2d{+}1),
\]
où $P(n) = \sum_{k=1}^{n} p_k$ est la somme des $n$~premiers premiers.
De manière équivalente,
\[
  \dim(d) = \sum_{k=1}^{d} (p'_k - 1) - (d - 1),
\]
où $p'_k = p_{k+1}$ est le $k$-ième premier impair.
\end{thmbox}

\begin{proof}[Esquisse de preuve]
La matrice d'observation à la profondeur~$d$ a $d$~blocs de lignes de
conditionnement, un par premier impair $p'_k \in \{3, 5, \ldots,
p'_d\}$. Après centrage, le bloc~$k$ contribue $p'_k - 1$ lignes
effectives. Le nombre total de lignes effectives est
$\sum_{k=1}^d (p'_k - 1)$.

Parmi ces lignes, exactement $d{-}1$ dépendances inter-blocs
existent~: chacun des $d{-}1$ «~anciens~» blocs (ceux déjà présents à
la profondeur $d{-}1$) partage une direction marginale avec l'espace
de profondeur précédente $V_{d-1}$. Le «~nouveau~» bloc (pour le
premier $p'_d$, entrant à la profondeur~$d$) n'introduit pas une
telle dépendance. Sous la position générale, ce sont les
\emph{seules} dépendances. Donc
$\dim(d) = \sum(p'_k - 1) - (d{-}1)$.
\end{proof}

\noindent
\textbf{Incrément.}
Une conséquence immédiate est
$\Delta\dim = \dim(d{+}1) - \dim(d) = p_{d+2} - 2$,
le nombre de transitions non triviales dans la matrice de transfert
$T_{p_{d+2}}$ de taille $(p_{d+2}{-}1) \times (p_{d+2}{-}1)$.
Le tableau~\ref{tab:PM_dim} vérifie la formule de dimension pour
$d = 2, \ldots, 11$.

\noindent
\textbf{Anticipation.} Le décalage de~2 (la couche à la
profondeur~$d$ dépend du premier $p_{d+2}$) coïncide avec
$\mathrm{PROF} = 2$, la profondeur du crible du théorème~T4.

\begin{table}[htbp]
\centering
\caption{Vérification de la dimension de couche (10/10).}
\label{tab:PM_dim}
\begin{tabular}{rccccc}
\toprule
$d$ & $P(d{+}1)$ & $2d{+}1$ & $\dim(d)$ prédit & observé & accord \\
\midrule
 2 &  10 &  5 &   5 &   5 & \checkmark \\
 3 &  17 &  7 &  10 &  10 & \checkmark \\
 4 &  28 &  9 &  19 &  19 & \checkmark \\
 5 &  41 & 11 &  30 &  30 & \checkmark \\
 6 &  58 & 13 &  45 &  45 & \checkmark \\
 7 &  77 & 15 &  62 &  62 & \checkmark \\
 8 & 100 & 17 &  83 &  83 & \checkmark \\
 9 & 129 & 19 & 110 & 110 & \checkmark \\
10 & 160 & 21 & 139 & 139 & \checkmark \\
11 & 197 & 23 & 174 & 174 & \checkmark \\
\bottomrule
\end{tabular}
\end{table}

\noindent
\textbf{Prédiction~:} $\dim(12) = P(13) - 25 = 238 - 25 = 213$.


% ============================================================
\subsection{Spectre d'activation}
\label{sec:PM_sigma}

\begin{thmbox}{Spectre de transversalité}
\label{thm:sigma}
Sous la conjecture~\ref{hyp:gen_pos} (AT-position générale), pour un
paquet de sondes de taille~$4$~:
\[
  \Sigma_k = \begin{cases}
    0 & \text{si } k < \mathrm{AT}, \\
    \displaystyle\binom{4}{k} & \text{si } k \ge \mathrm{AT}.
  \end{cases}
\]
\end{thmbox}

\begin{proof}
Soit $W = \operatorname{span}(u_1, \ldots, u_4)$ l'espace transverse
(les projections des 4~sondes sur $V^{\perp}$), avec
$\dim W = \mathrm{AT}$.

\emph{Cas $k < \mathrm{AT}$~:} un sous-ensemble de $k < \mathrm{AT}$
sondes engendre un espace de dimension $\le k < \mathrm{AT}$ et ne
peut couvrir~$W$.

\emph{Cas $k \ge \mathrm{AT}$~:} tout sous-ensemble de taille~$k$
contient un sous-sous-ensemble de taille~AT qui, par position
générale, a un rang exactement~AT et engendre~$W$. Donc tout
sous-ensemble de ce type détecte le DDL, donnant
$\Sigma_k = \binom{4}{k}$.
\end{proof}

\noindent
Vérification~: 13/13 exact (tableau~\ref{tab:PM_shells}).


% ============================================================
\subsection{Seuil d'activation~: trois régimes}
\label{sec:PM_AT}

Le seuil d'activation est gouverné par
$|B| = |B(p)| = \pi(p) - 4$,
le nombre de niveaux de conditionnement d'écho (premiers entre 7
et~$p$, exclus).

\subsubsection{Partie~A~: réinitialisation perturbative TCR (cycles $\ge 2$)}

\begin{thmbox}{Seuil d'activation A}
\label{thm:AT_A}
Si $|B(p)| > \varphi(30) = 8$ (de manière équivalente~: il existe
$q < p$ tel que $q \equiv p \pmod{30}$), alors $\mathrm{AT}(p) = 1$.
\end{thmbox}

\begin{proof}
Pour $|B| \ge 9$ (c'est-à-dire $p \ge 41$), il existe un prédécesseur
$q < p$ avec $q \equiv p \pmod{30}$, c.-à-d.\ $q \equiv p \pmod{2}$,
$q \equiv p \pmod{3}$, $q \equiv p \pmod{5}$. La composante TCR
mod~30 de $\eta_{p,c}$ est déjà capturée par la sonde existante
$\eta_{q,c'}$ dans la base. La composante complémentaire
(distinguant $p$ de~$q$ à l'intérieur de leur canal mod-30 partagé)
est de rang~1. Une seule sonde $\eta_{p,c}$ suffit à capturer cette
correction de rang~1. Par L1, la profondeur future ajoute exactement
un DDL, donc $\mathrm{AT} = 1$.
\end{proof}

\subsubsection{Partie~B~: régime d'écho ($|B| \le 2$)}

\begin{thmbox}{Seuil d'activation B}
\label{thm:AT_B}
Pour $p \in \{11, 13\}$ (premiers d'écho, $|B| \le 2$)~:
$\mathrm{AT}(p) = 1$.
\end{thmbox}

\begin{proof}
Les premiers 11 et 13 sont les premiers premiers d'écho
($\gamp{11} = 0{,}426$, $\gamp{13} = 0{,}356$, tous deux $< 1/2$). À
leur profondeur, la base ne contient aucune sonde pour les résidus
mod~11 (resp.\ 13). Tout l'espace de sondes $\eta_{p,c}$ est nouveau.
Une seule sonde capture l'unique nouveau DDL garanti par~L1.
\end{proof}

\subsubsection{Partie~C~: mécanisme de rang fantôme (cycle~1)}

Le c\oe ur du chapitre. On établit d'abord le phénomène de
\emph{rang fantôme}.

\begin{definition}[Rang fantôme et rang structurel]
\label{def:phantom_struct}
Soit $F_0$ la famille de base ($n_{\mathrm{base}}$ sondes) et
$F_k = F_0 \cup \{\eta_1, \ldots, \eta_k\}$ la famille étendue avec
$k$ sondes-cibles. On définit~:
\begin{align}
  \mathrm{fantome}(k) &= R(F_k, d) - R(F_0, d), \label{eq:phantom}\\
  \mathrm{struct}(k) &= R(F_k, d{+}1) - R(F_k, d), \label{eq:struct}
\end{align}
où $d = \mathrm{prof\_base}$ et $d{+}1$ est la profondeur future.
Alors $\mathrm{AT} = \min\{k : \mathrm{struct}(k) > 0\}$.
\end{definition}

\begin{idbox}{C1 --- Identité de partition}
\label{id:C1}
\mID
$\mathrm{fantome}(1) + \mathrm{struct}(1) = n_{\mathrm{base}}$.
\end{idbox}

\begin{proof}
La famille de base observe toutes les profondeurs jusqu'à~$d$ et ne
contient aucune sonde pour le premier-cible~$p$. Donc le rang
structurel de base seul est nul~:
$\dim_{\mathrm{base}} = R(F_0, d{+}1) - R(F_0, d) = 0$. Alors~:
\begin{align*}
  \mathrm{fantome}(1) + \mathrm{struct}(1)
  &= [R(F_1, d) - R(F_0, d)] + [R(F_1, d{+}1) - R(F_1, d)] \\
  &= R(F_1, d{+}1) - R(F_0, d) \\
  &= [R(F_1, d{+}1) - R(F_0, d{+}1)] + \dim_{\mathrm{base}} \\
  &= R(F_1, d{+}1) - R(F_0, d{+}1).
\end{align*}
Sous position générale, les $n_{\mathrm{base}}$ termes croisés
$(\eta_{\mathrm{cible}} \times \text{base}_j)$ à la profondeur
$d{+}1$ sont linéairement indépendants, donc
$R(F_1, d{+}1) - R(F_0, d{+}1) = n_{\mathrm{base}}$.
\end{proof}

\noindent
Le tableau~\ref{tab:C1} confirme l'identité de partition pour les
six strates testées.

\begin{table}[htbp]
\centering
\caption{Identité de partition C1~: vérification exacte (6/6).}
\label{tab:C1}
\begin{tabular}{rrrrrr}
\toprule
Strate & $n_{\mathrm{base}}$ & fantome(1) & struct(1) & ph + st & accord \\
\midrule
11 &  6 &  4 & 2 &  6 & \checkmark \\
13 & 10 &  8 & 2 & 10 & \checkmark \\
17 & 14 & 14 & 0 & 14 & \checkmark \\
19 & 18 & 18 & 0 & 18 & \checkmark \\
23 & 22 & 22 & 0 & 22 & \checkmark \\
29 & 26 & 26 & 0 & 26 & \checkmark \\
\bottomrule
\end{tabular}
\end{table}

\begin{thmbox}{C2 --- Saturation du fantôme}
\label{thm:C2}
Sous la conjecture~\ref{hyp:gen_pos}, pour $|B| \ge 3$~:
$\mathrm{fantome}(1) = n_{\mathrm{base}}$, donc
$\mathrm{struct}(1) = 0$.
\end{thmbox}

\begin{proof}
Lorsque $p$ ne divise pas $M_d = \prod_{k=1}^{d} p_k$, la
distribution de $\operatorname{surv}(d) \bmod p$ est non uniforme.
La décomposition TCR de la matrice d'observation révèle
$N_{\mathrm{spatial}} = 3$ canaux actifs (mod~3, mod~5, mod~7). La
non-uniformité crée des corrélations entre $\eta_{p,c}$ et les
sondes de base à la profondeur~$d$.

Pour $|B| \le 2$~: seulement 1--2 niveaux de conditionnement d'écho~;
les 3~canaux TCR ne sont pas tous saturés. Deux directions
structurelles survivent~: $\mathrm{struct}(1) = 2 > 0$.

Pour $|B| \ge 3$~: les $|B|$ niveaux d'écho fournissent assez de
non-uniformité pour saturer les 3~canaux TCR. Tous les
$n_{\mathrm{base}}$ termes croisés sont absorbés par le fantôme~:
$\mathrm{struct}(1) = 0$.

Le seuil $|B| = 3$ est un comptage dimensionnel~: il faut au moins
$N_{\mathrm{spatial}} = 3$ sources indépendantes de non-uniformité
pour bloquer les 3~canaux TCR.
\end{proof}

\begin{thmbox}{Seuil d'activation C}
\label{thm:AT_C}
Sous la conjecture~\ref{hyp:gen_pos}, pour les strates non-écho du
cycle~1 ($3 \le |B| \le 8$)~:
\[
  \mathrm{AT} = \left\lceil \frac{|B| - 2}{N_{\mathrm{spatial}}} \right\rceil + 1,
  \qquad N_{\mathrm{spatial}} = |\{3, 5, 7\}| = 3.
\]
\end{thmbox}

\begin{proof}
Par le théorème~\ref{thm:C2}, $\mathrm{struct}(1) = 0$ pour $|B| \ge
3$. L'argument procède en trois étapes.

\emph{Étape~1 (partition des canaux TCR).} La matrice d'observation
se décompose en $N_{\mathrm{spatial}} = 3$ canaux (mod~3, mod~5,
mod~7) via l'isomorphisme TCR $\Z/M_d^*\Z \cong \prod \Z/p_i^*\Z$.
Le fantôme sature chaque canal indépendamment.

\emph{Étape~2 (saturation par canal).} Chaque sonde-cible crée des
termes croisés dans les 3~canaux TCR. Sous position générale, chaque
sonde contribue jusqu'à $N_{\mathrm{spatial}} = 3$ directions
inter-canaux brisant le fantôme. Le surplus de fantôme à surmonter
est $|B| - 2$ (les $|B|$ niveaux d'écho moins les 2 niveaux
«~libres~» qui ne bloquent pas tous les canaux).

\emph{Étape~3 (comptage).} Pour $\mathrm{struct}(k) > 0$, les
$k$~sondes doivent collectivement épuiser le fantôme~:
\[
  k \cdot N_{\mathrm{spatial}} \ge |B| - 2
  \;\Longrightarrow\;
  \mathrm{AT} = \left\lceil\frac{|B|-2}{3}\right\rceil + 1.
\]
Le $+1$ rend compte de la sonde de référence requise même sans
fantôme.
\end{proof}

\noindent
Le tableau~\ref{tab:AT_verify} vérifie la formule à travers les trois
régimes pour 13~strates.

\begin{table}[htbp]
\centering
\caption{Formule AT~: vérification complète (13/13).}
\label{tab:AT_verify}
\begin{tabular}{rrcccl}
\toprule
Strate & $|B|$ & Régime & AT (obs.) & AT (formule) & Accord \\
\midrule
11 & 1 & écho     & 1 & 1 & \checkmark \\
13 & 2 & écho     & 1 & 1 & \checkmark \\
17 & 3 & constr.  & 2 & $\lceil 1/3\rceil + 1 = 2$ & \checkmark \\
19 & 4 & constr.  & 2 & $\lceil 2/3\rceil + 1 = 2$ & \checkmark \\
23 & 5 & constr.  & 2 & $\lceil 3/3\rceil + 1 = 2$ & \checkmark \\
29 & 6 & constr.  & 3 & $\lceil 4/3\rceil + 1 = 3$ & \checkmark \\
31 & 7 & constr.  & 3 & $\lceil 5/3\rceil + 1 = 3$ & \checkmark \\
37 & 8 & constr.  & 3 & $\lceil 6/3\rceil + 1 = 3$ & \checkmark \\
41 & 9 & asympt.  & 1 & 1 & \checkmark \\
43 & 10 & asympt. & 1 & 1 & \checkmark \\
47 & 11 & asympt. & 1 & 1 & \checkmark \\
53 & 12 & asympt. & 1 & 1 & \checkmark \\
59 & 13 & asympt. & 1 & 1 & \checkmark \\
\bottomrule
\end{tabular}
\end{table}


% ============================================================
\subsection{Rang fantôme~: un nouveau phénomène arithmétique}
\label{sec:PM_phantom}

\subsubsection{Origine}
\mTHM

Le rang fantôme provient d'une cause purement arithmétique~: lorsque
$p \nmid M_d$, l'ensemble $\operatorname{surv}(d) \bmod p$ n'est
\emph{pas} équidistribué. Les classes de résidus portent des biais
non triviaux qui créent des corrélations parasites entre la
sonde-cible $\eta_{p,c}$ et les sondes de base. Ces corrélations
gonflent le rang à la profondeur de base~$d$ sans contribuer aucune
information structurelle à la profondeur future $d{+}1$.

\subsubsection{Loi de conservation}

L'identité~\ref{id:C1} est une \emph{loi de conservation}~:
l'information totale créée en ajoutant une sonde-cible
($= n_{\mathrm{base}}$ nouvelles colonnes de termes croisés) se
sépare en composantes fantôme et structurelle dont la somme est
constante. Pour les strates d'écho ($|B| \le 2$)~:
$\mathrm{struct}(1) = 2$,
$\mathrm{fantome}(1) = n_{\mathrm{base}} - 2$. Pour les strates
non-écho ($|B| \ge 3$)~: $\mathrm{struct}(1) = 0$,
$\mathrm{fantome}(1) = n_{\mathrm{base}}$.

\subsubsection{Le fantôme comme prix de la mémoire arithmétique}

Le rang fantôme est le \emph{prix} de la mémoire arithmétique~: la
non-uniformité de $\operatorname{surv}(d) \bmod p$ enregistre quels
premiers ont déjà été criblés. Cette mémoire gonfle le rang à la
profondeur~$d$, masquant le signal structurel. Le seuil
d'activation AT mesure combien de sondes sont nécessaires pour
«~percer~» la barrière fantôme.

La constante $N_{\mathrm{spatial}} = 3$ dans la formule AT est le
nombre de canaux TCR actifs. C'est la même quantité qui détermine
la dimensionnalité spatiale dans PT (article
compagnon~\cite{companion_M3}). Le mécanisme de rang fantôme fournit
ainsi une interprétation \emph{informationnelle} concrète de
$N_{\mathrm{spatial}} = 3$~: c'est la constante d'obstruction à la
détection structurelle à travers le bruit arithmétique.

\subsubsection{Contraste avec les codes linéaires}

Dans la seconde instanciation (codes correcteurs d'erreurs sur
$\mathbb{F}_2$~; section~\ref{sec:PM_codes}), le rang fantôme est
\emph{identiquement nul}~: $\mathrm{AT} = 1$ universellement. Cela
provient de ce que les codes $\mathbb{F}_2$-linéaires ont des
distributions de syndromes uniformes~; aucune mémoire arithmétique
ne s'accumule. Le fantôme est spécifique au crible, pas universel.


% ============================================================
\subsection{Seconde instanciation~: codes correcteurs d'erreurs}
\label{sec:PM_codes}

Le cadre MP a été testé sur un second domaine~: les codes linéaires
correcteurs d'erreurs sur $\mathbb{F}_2$.

\subsubsection{Correspondance structurelle}

\begin{center}
\begin{tabular}{ll}
\toprule
\textbf{Crible} & \textbf{Code linéaire} \\
\midrule
Entiers $1, \ldots, N$ & Mots de code $\mathbb{F}_2^n$ \\
Premier $p_k$ & Contrôle de parité $h_k$ \\
«~Non divisible par $p$~» & $h_k \cdot x = 0 \bmod 2$ \\
Survivants $d$-rugueux & $C_d = \ker(H_d)$ \\
Sonde $\eta(p, c)$ & Sonde $\psi_j = x_j - 1/2$ \\
Strate mod~30 & Coset cyclotomique (BCH) \\
\bottomrule
\end{tabular}
\end{center}

\subsubsection{Résultats}

Sept codes ont été testés (Répétition $[5,1,5]$, Hamming $[7,4,3]$,
Hamming étendu $[8,4,4]$, systématique $[6,3]$, BCH $[15,5,7]$,
Reed--Muller $[8,4,4]$, Golay $[23,12,7]$). Dans tous les cas~:
\begin{itemize}[nosep]
  \item $\mathrm{AT} = 1$ universellement (pas de rang fantôme).
  \item L1 tient conditionnellement (lorsque les sondes s'étendent
    avec les contraintes).
  \item La matrice d'observation, la matrice historique et le rang
    cumulatif sont bien définis et calculables.
\end{itemize}

\subsubsection{Classification~: universel vs.\ spécifique}

Le tableau~\ref{tab:PM_classification} classifie chaque résultat MP
par sa portée de validité.

\begin{table}[htbp]
\centering
\caption{Résultats MP universels vs.\ spécifiques au crible.}
\label{tab:PM_classification}
\begin{tabular}{lll}
\toprule
\textbf{Résultat} & \textbf{Statut} & \textbf{Portée} \\
\midrule
Cadre MP (obs., hist., rang) & universel & tout $(F, C)$ \\
T1 (transitions interdites) & quasi-universel & tout crible de pas~3 \\
$\spt = 1/2$ & quasi-universel & tout crible de pas~3 \\
L1 (incrément unité) & spécifique & Ératosthène (TCR requis) \\
AT $= 1$ & universel (linéaire) & codes $\mathbb{F}_2$ \\
AT $= 1$ & spécifique (fantôme) & crible, cycles $\ge 2$ \\
$\dim(d) = P(d{+}1) - (2d{+}1)$ & spécifique & crible seul (TCR requis) \\
$\Sigma_k = \binom{4}{k}$ & spécifique & crible (4 composantes TCR) \\
Rang fantôme & spécifique & mémoire arithmétique \\
Cosets $\sim$ strates & structurel & cosets BCH $\sim$ strates mod~30 \\
\bottomrule
\end{tabular}
\end{table}


% ============================================================
\subsection{Troisième instanciation~: nombres chanceux}
\label{sec:PM_lucky}

Le crible des nombres chanceux est un crible \emph{positionnel}~:
partant des entiers impairs (pas~0 $= 2$), il retire chaque
$L_k$-ième survivant à l'étape~$k$, où $L_k$ est le $k$-ième élément
restant. Puisqu'Ératosthène est l'unique crible primitif (T6), le
crible chanceux est un \emph{crible sur le crible}~: un raffinement
positionnel du substrat arithmétique.

\subsubsection{Cadre}

Valeurs du crible chanceux~: $2, 3, 7, 9, 13, 15, 21, 25, 31, 33, 37,
\ldots$ (seulement 47\% sont premiers~--- le reste est composé comme
9, 15, 21, 25). À chaque profondeur~$d$, les survivants et leurs
sauts sont analysés avec les mêmes sondes MP (indicateurs
modulaires mod~3, 5, 7) et structure de conditionnement que pour
Ératosthène.

\subsubsection{T0 et $\spt = 1/2$~: universels}

Le tableau~\ref{tab:PM_lucky} compare les invariants MP pour les
deux cribles à $N = 10^5$.

\begin{table}[htbp]
\centering
\caption{Comparaison MP~: Ératosthène vs.\ nombres chanceux ($N = 10^5$).}
\label{tab:PM_lucky}
\begin{tabular}{lll}
\toprule
\textbf{Propriété} & \textbf{Ératosthène} & \textbf{Chanceux} \\
\midrule
T1 (transitions interdites mod~3) & OUI (exact) & OUI (exact) \\
$\spt = n_1/(n_1 + n_2)$ & $0{,}49983$ & $0{,}50000$ (exact) \\
$|\spt - 1/2|$ & $1{,}7 \times 10^{-4}$ & $0$ \\
$T_{00}$ (auto-trans.\ classe~0) & $0{,}362$ & $0{,}270$ \\
L1 (incrément unité) & OUI (13/13) & NON (alterné 2--1) \\
Saturation de rang & $\infty$ (TCR) & 12 (fini) \\
Type de contrainte & Divisibilité & Positionnel \\
Factorisation TCR & OUI & NON \\
\bottomrule
\end{tabular}
\end{table}

\noindent
Le paramètre de symétrie $\spt = 1/2$ tient \emph{exactement} pour
les nombres chanceux ($n_1 = n_2 = 5725$ à $N = 10^5$), encore plus
précisément que pour les premiers. Le mécanisme est le même~: T0
interdit les transitions non nulles intra-classe, forçant
l'involution $1 \leftrightarrow 2$ et donc $n_1 = n_2$.

\begin{remarque}
T0 pour les nombres chanceux ne vient \emph{pas} de la divisibilité.
Il vient du pas positionnel $L_2 = 3$ (retirant chaque 3-ième
élément), qui crée une contrainte mod-3 équivalente. Tout crible
dont le second pas non trivial est~3 hérite de T0. Puisqu'Ératosthène
commence avec les pas 2 et 3, et que les chanceux commencent avec
les pas 2 et 3 (positionnellement), tous deux héritent de T0 et
de $\spt = 1/2$.
\end{remarque}

\subsubsection{L1 et structure de rang~: spécifique à Ératosthène}

La structure de rang chanceuse est qualitativement différente~:
\[
  \dim_{\mathrm{Chanceux}}(d) = \begin{cases}
    2 & d \text{ impair, } d \le 7, \\
    1 & d \text{ pair, } d \le 6, \\
    0 & d \ge 8.
  \end{cases}
\]
Les rangs saturent à $12 = \sum_{m \in \{3,5,7\}} (m - 1)$, le
maximum possible avec des sondes mod $\{3, 5, 7\}$. Sans termes
croisés TCR, le crible positionnel ne peut créer d'observables
irréductibles aux modules composés ($15, 21, 35, 105$), donc la
croissance du rang est bornée par l'alphabet des sondes.

À l'inverse, le crible d'Ératosthène produit des observables
TCR-irréductibles à chaque profondeur (lemme~\ref{lem:CRT_irred}),
donnant une croissance non bornée du rang et des incréments unités
exacts (L1).

\subsubsection{Mise à jour de la classification}

L'instanciation chanceuse affine la classification
universel/spécifique~:

\begin{center}
\begin{tabular}{llll}
\toprule
\textbf{Résultat} & \textbf{Codes ($\mathbb{F}_2$)} & \textbf{Chanceux} & \textbf{Ératosthène} \\
\midrule
$\spt = 1/2$ & N/A & OUI (exact) & OUI \\
T0 & N/A & OUI & OUI \\
L1 & conditionnel & NON & OUI \\
$\dim(d) = P(d{+}1) - (2d{+}1)$ & NON & NON & OUI \\
$\Sigma_k = \binom{4}{k}$ & N/A & N/A & OUI \\
Formule AT (3 régimes) & $\mathrm{AT} = 1$ & N/A & OUI \\
Rang fantôme & absent & différent & OUI \\
\bottomrule
\end{tabular}
\end{center}

\noindent
La hiérarchie est claire~: $\spt = 1/2$ et T0 sont \emph{universels}
(tout crible de pas~3)~; L1 et la formule de dimension sont
\emph{spécifiques} à Ératosthène (requièrent TCR)~; le mécanisme de
rang fantôme est \emph{spécifique} à la mémoire arithmétique
(requiert $\operatorname{surv}(d) \bmod p$ non uniforme).

\subsubsection{Convergence de l'holonomie chanceuse}
\label{sec:PM_lucky_holonomy}

Le crible chanceux partage T0 et $\spt = 1/2$ avec Ératosthène mais
manque de multiplicativité. Néanmoins, les angles d'holonomie
convergent asymptotiquement~:

\begin{itemize}[nosep]
\item Pour $L_k \ge 31$~:
  $\sinp{\mathrm{Chanceux}} / \sinp{\mathrm{PT}} \to 1{,}00$.
\item $\gamp{\mathrm{Chanceux}}$ oscille plutôt que de converger
  monotonement~: 506/991 incréments sont positifs (50{,}9\%,
  compatible avec une marche aléatoire).
\item La somme cumulative $\sum L_k = 4\,057\,742 \gg \mustar
  \approx 12$ (point fixe chanceux), confirmant l'absence de
  convergence finie.
\end{itemize}

Cette convergence est \emph{asymptotique}, non structurelle~: les
nombres chanceux approchent les valeurs PT par le haut sans le
mécanisme TCR qui rend Ératosthène exact. Le crible chanceux est une
ombre d'Ératosthène, partageant sa symétrie ($\spt$) mais pas sa
géométrie ($\sinp{p}$).

\begin{remarque}
Le comportement oscillatoire des angles d'holonomie chanceux, avec
incréments quasi-symétriques (50{,}9\% positifs), est
caractéristique d'un crible positionnel~: en l'absence de
multiplicativité, chaque étape d'élimination crée une perturbation
en marche aléatoire autour de la valeur asymptotique. Le taux de
convergence
$|1 - \sinp{\mathrm{Chanceux}} / \sinp{\mathrm{PT}}| \sim L_k^{-1/2}$
est cohérent avec cette image et contraste fortement avec la
convergence exponentielle de $\sinp{p}$ sous Ératosthène.
\end{remarque}


% ============================================================
\subsection{Le lien MP--PT}
\label{sec:PM_PT}

\begin{corollaire}[Conjecture MP--PT, prouvée]
\label{cor:PMPT}
$\mathrm{AT}(p) = 1$ pour tout premier $p \ge 41$.
\end{corollaire}

\begin{proof}
Corollaire direct du théorème~\ref{thm:AT_A}~: pour $p \ge 41$,
$|B| \ge 9 > \varphi(30) = 8$, donc la réinitialisation perturbative
s'applique.
\end{proof}

\subsubsection{Conséquences}

\begin{enumerate}[nosep]
  \item \textbf{Complétude du MS.} L'escalator d'activation (cycle~1,
    strates 11--37) est un \emph{transitoire}~: la complexité
    structurelle du crible est épuisée en 8~strates. Toutes les
    strates ultérieures sont des corrections perturbatives
    ($\mathrm{AT} = 1$). Les 43~observables du MS sont le contenu
    \emph{complet} de ce transitoire.

  \item \textbf{Transitoire = physique.} La complexité du monde
    physique (3~générations, 3~dimensions spatiales, 41~observables)
    est entièrement contenue dans le premier cycle primoriel. Le
    régime permanent ($\mathrm{AT} = 1$) est le silence structurel.

  \item \textbf{Pas de nouvelle physique.} La structure du crible est
    épuisée au cycle~1. Les premiers d'écho ($p \ge 11$) n'ajoutent
    que des corrections perturbatives (PV d'écho), jamais une
    nouvelle structure qualitative. Prédiction~: aucune nouvelle
    particule, force ou symétrie au-delà du Modèle Standard.
\end{enumerate}

\subsubsection{Correspondances d'invariants}

\begin{center}
\begin{tabular}{ll}
\toprule
\textbf{Invariant MP} & \textbf{Invariant PT} \\
\midrule
Profondeur future (PF) & Position dans la cascade vers $\mustar$ \\
Seuil d'activation (AT) & «~Résistance~» du premier~$p$ \\
Spectre $\Sigma$ & Géométrie d'holonomie du paquet de sondes \\
Dimension de couche & Liée à $\varphi(p)$ d'Euler \\
Cycle primoriel & $30 = 2\mustar$ (primoriel actif) \\
Décalage de stabilisation $= 2$ & $\mathrm{PROF} = 2$ \\
$N_{\mathrm{spatial}} = 3$ dans AT & $|\{3, 5, 7\}| = 3$ (premiers actifs) \\
\bottomrule
\end{tabular}
\end{center}


% ============================================================
\subsection{Hiérarchie des invariants de persistance}
\label{sec:PM_hierarchy}

Le cadre MP révèle une hiérarchie à quatre niveaux d'invariants,
ordonnés par universalité (tableau~\ref{tab:PM_hierarchy}).

\begin{table}[htbp]
\centering
\caption{Quatre niveaux d'invariants de persistance.}
\label{tab:PM_hierarchy}
\begin{tabular}{clll}
\toprule
\textbf{Niveau} & \textbf{Nom} & \textbf{Invariant} & \textbf{Portée} \\
\midrule
0 & Universel & $\DKL(P \| U)$, $H_{\max} = \DKL + H$ & Tout système T0 \\
1 & Type & $\spt = 1/2$, T0 transitions interdites & Cribles de type~I \\
2 & Cascade & $\mustar = 15$, $\{3,5,7\}$ actifs & TCR-compatibles \\
3 & Crible & $\sinp{p}$, $\gamp{p}$, $\alphaEM$ & Ératosthène seul \\
\bottomrule
\end{tabular}
\end{table}

\paragraph{Niveau~0~: universel.}
L'identité $\log_2 3 = \DKL + H$ tient \emph{exactement}
($< 10^{-15}$) pour \emph{tous} les systèmes T0~: le crible
d'Ératosthène, le crible chanceux, le système T0 des protéines, les
sauts de marche aléatoire, et les suites de classes alternées. C'est
le contenu du TFS (Théorème Fondamental des Sauts, article
compagnon~\cite{companion_M1})~: la séparation de l'entropie maximale
$H_{\max}$ en divergence $\DKL$ et entropie résiduelle~$H$ est une
conséquence universelle de la contrainte T0, sans dépendance au
crible particulier ni à l'origine des sauts.

\paragraph{Niveau~1~: type.}
Le paramètre de symétrie $\spt = 1/2$ et les transitions interdites
T0 tiennent à la fois pour Ératosthène et les nombres chanceux
(section~\ref{sec:PM_lucky}). Le mécanisme est le même~: tout crible
dont le second pas non trivial est~3 hérite de la contrainte mod-3
qui force $n_1 = n_2$. Les invariants de niveau~1 sont partagés par
tous les \emph{cribles de type~I} (pas~3 présent).

\paragraph{Niveau~2~: cascade.}
Le point fixe $\mustar = 15$ et l'ensemble des premiers actifs
$\{3,5,7\}$ requièrent l'indépendance multiplicative (structure TCR).
Le crible chanceux, positionnel et non multiplicatif, n'atteint
\emph{pas} $\mustar = 15$~: son rang sature à~12 et sa convergence
est asymptotique, non structurelle
(section~\ref{sec:PM_lucky_holonomy}).

\paragraph{Niveau~3~: crible.}
Les angles d'holonomie $\sinp{p}$ sont définis par la géométrie
spécifique du crible d'Ératosthène au point fixe $\mustar = 15$. Les
nombres chanceux ne possèdent pas ces angles~; les codes correcteurs
d'erreurs ne les possèdent pas~; les protéines ne les approchent que
statistiquement. Le niveau~3 est la \emph{signature géométrique}
d'Ératosthène.

\paragraph{Théorème de sélection.}
\mTHM
Parmi cinq systèmes candidats~--- Ératosthène, chanceux, protéines,
marche aléatoire, alterné~---, \emph{seul Ératosthène} satisfait les
quatre conditions simultanément~:
\begin{enumerate}[nosep]
  \item Type~I (T0 avec pas~3),
  \item $\spt = 1/2$ exact,
  \item Multiplicativité TCR,
  \item Cascade infinie (croissance non bornée du rang).
\end{enumerate}
Chanceux échoue à la condition~3 (pas de TCR). Les protéines
échouent à la condition~2 ($\spt$ approximatif, pas exact). La
marche aléatoire et les suites alternées échouent à la condition~1
(pas de contrainte T0 issue du pas~3).

\paragraph{Spécificité de $\sinp{p}$.}
Six approches (étiquetées A--F) ont été testées pour dériver
$\sinp{p}$ des valeurs propres empiriques de la matrice de
transfert. Toutes ont donné des résultats négatifs~: écart de 12\%,
10\%, 8{,}5\% pour $p = 3, 5, 7$ respectivement. Les angles
d'holonomie $\sinp{p}$ sont définis par la théorie ($\qplus$ au
point fixe $\mustar = 15$), et non par la matrice $T$ empirique. La
matrice de transfert encode les \emph{corrélations} (dynamique),
tandis que $\sinp{p}$ encode la \emph{géométrie} (holonomie)~---
aspects complémentaires du même crible, non dérivables l'un de
l'autre.

\begin{remarque}
Les profondeurs d'observation de l'article
compagnon~\cite{companion_chemistry} instancient cette hiérarchie~:
chaque domaine~--- chimie, nucléaire, hadronique~--- mesure le même
$\DKL$ à une profondeur d'observation~$d$ différente. Le tableau
ci-dessus organise ces domaines verticalement, de l'universel
(niveau~0, partagé par tous) au spécifique (niveau~3, unique à
Ératosthène).
\end{remarque}


% ============================================================
\subsection{La structure auto-référentielle des MP}
\label{sec:PM_self}

\subsubsection{Les MP ne sont pas un outil externe}

Les instruments mathématiques standard pour prédire des invariants
structurels~--- le théorème d'indice d'Atiyah--Singer, la capacité de
Shannon, l'homologie persistante~--- opèrent tous \emph{sur} une
structure donnée extérieurement~: une variété, un canal, un complexe
simplicial. Les MP occupent une position catégoriquement
différente~: \emph{le crible prédit son propre contenu en
information}.

La formule de dimension de couche $\dim(d) = P(d{+}1) - (2d{+}1)$
utilise la fonction de somme des premiers $P(n) = \sum_{k=1}^n p_k$~---
construite à partir de la sortie du crible~--- pour prédire la
structure de rang de la matrice d'observation~--- elle-même construite
à partir des contraintes du crible. La formule AT utilise
$|B| = \pi(p) - 4$, là encore une sortie du crible, pour prédire le
coût d'activation de la prochaine strate. C'est une auto-référence
sans circularité~: le \emph{niveau de description} (rangs, DDL,
spectres) est distinct du \emph{niveau calculé} (survivants, sauts,
résidus).

\begin{remarque}
Dans le langage de \cref{sec:math_structures}, MP n'est pas une
\emph{application} des structures vers des invariants~; c'est une
\emph{propriété de point fixe} de la géométrie d'information du
crible lui-même. Le crible n'a pas besoin d'observateur~--- il
s'observe lui-même, et l'observation sature à la 13$^{\rm e}$~strate.
\end{remarque}

\subsubsection{Dissolution de la dichotomie discret--continu}

Le cadre MP dissout la frontière apparente entre les mathématiques
discrètes et continues. À chaque niveau, les mêmes objets MP ont
deux faces~:

\begin{center}
\begin{tabular}{lll}
\toprule
\textbf{Objet MP} & \textbf{Face discrète} & \textbf{Face continue} \\
\midrule
Matrice d'observation & Rang entier & Métrique de Fisher \\
Profondeur de contrainte $d$ & Indice de crible $k$ & Paramètre d'échelle $\mu$ \\
AT $= 1$ (asymptotique) & Comptage de sondes & Convergence $Q \to 0$ \\
Rang fantôme & Défaut de rang (entier) & Courbure d'information \\
$\binom{4}{k}$ & Combinatoire finie & Limite binomiale \\
\bottomrule
\end{tabular}
\end{center}

\noindent
Ce ne sont pas des approximations l'une de l'autre. L'angle
d'holonomie $\sinp{p} = \delta_p(2 - \delta_p)$ relie $\Z/p\Z$ (un
groupe fini) à $S^1$ (le cercle continu) via l'identité
$\zeta(2) = \pi^2/6 = \prod_p (1 - 1/p^2)^{-1}$~: le nombre $\pi$
\emph{émerge} du crible. La variable complexe $w_p$
(\cref{sec:complex}) étend les matrices de transfert discrètes
$T_p$ en systèmes dynamiques holomorphes. Travailler sur le crible,
c'est simultanément travailler sur les entiers, les réels et le plan
complexe.

\begin{remarque}
La dissolution de la limite continue (R50) a montré que la métrique
de Fisher \emph{est} la limite continue~--- non une approximation,
mais une identité. MP fournit la contrepartie structurelle~: la
hiérarchie de rangs de la matrice d'observation \emph{est} le
squelette discret de la même géométrie. Discret et continu ne sont
pas deux mondes~; ce sont deux descriptions de la même structure
arithmétique, et MP est l'endroit où cette identité devient
explicite calculatoirement.
\end{remarque}


% ============================================================
\subsection{Dissolution de l'extérieur}
\label{sec:PM_exterior}

Quatre questions traditionnellement considérées comme «~externes~» aux
mathématiques se dissolvent dans le cadre MP~:

\paragraph{Q1~: implémentation.}
«~Qui exécute le crible~?~»~--- Personne. Le crible n'est pas un
processus~; c'est une cascade de contraintes. Le crible
d'Ératosthène existe comme objet mathématique, et non comme
algorithme nécessitant un ordinateur. Le cadre MP montre que son
contenu en information est auto-déterminé
(section~\ref{sec:PM_self}).

\paragraph{Q2~: réflexivité.}
«~Le crible peut-il décrire sa propre complexité~?~»~--- Oui. La
formule de dimension de couche $\dim(d) = P(d{+}1) - (2d{+}1)$
utilise la fonction de somme des premiers~--- une sortie du
crible~--- pour prédire la structure de rang de la matrice
d'observation~--- une entrée du crible. Cette auto-référence n'est
pas circulaire~: le niveau de description diffère du niveau calculé.

\paragraph{Q3~: conscience.}
«~Le crible requiert-il un observateur~?~»~--- Non. Le mécanisme de
rang fantôme MP opère identiquement, que quelqu'un le calcule ou
non. Le seuil d'activation AT est une propriété de l'algèbre TCR, et
non d'un agent observateur.

\paragraph{Q4~: choix ontologique.}
«~Le crible est-il physiquement réel ou une construction
mathématique~?~»~--- Cette question est dissoute, et non résolue. Les
trois lectures ontologiques de la monographie compagnon (coïncidence,
correspondance, réalité) sont toutes compatibles avec le contenu
mathématique. PT ne requiert aucun engagement métaphysique~; elle
requiert uniquement que l'algèbre du crible soit correctement
calculée.

\begin{remarque}
Ces dissolutions sont des \emph{conséquences} de MP, et non des
opinions philosophiques. Q1 est dissoute par la structure de rang
auto-référentielle (section~\ref{sec:PM_self})~; Q2 par la formule de
dimension de couche (théorème~\ref{thm:dim_formula})~; Q3 par
l'indépendance vis-à-vis de l'observateur du rang fantôme~; Q4 par
la multiplicité des lectures ontologiques compatibles. Dans chaque
cas, la question présuppose une distinction (algorithme vs.\ objet,
descripteur vs.\ décrit, observateur vs.\ observé, réel vs.\
abstrait) que MP efface structurellement.
\end{remarque}


% ============================================================
\subsection{Résumé et architecture}
\label{sec:PM_summary}

Le cadre MP fournit 7~résultats prouvés à partir d'une unique
fondation algébrique (algèbre des indicateurs + TCR)~:

\begin{center}
\begin{tikzpicture}[
  node distance=1.2cm and 2cm,
  every node/.style={draw, rounded corners, minimum width=2.8cm,
    minimum height=0.7cm, font=\small, align=center},
  thm/.style={fill=thmblue!10, draw=thmblue},
  idstyle/.style={fill=idgreen!10, draw=idgreen},
  arr/.style={-{Stealth}, thick}
]
\node[thm] (L1) {L1\\$\Delta\mathrm{PF} = 1$\\{\scriptsize\THMTAG}};
\node[thm, below left=of L1] (dim) {$\dim(d)$\\$= P(d{+}1){-}(2d{+}1)$\\{\scriptsize\THMTAG*}};
\node[thm, below=of L1] (sig) {$\Sigma_k = \binom{4}{k}$\\{\scriptsize\THMTAG*}};
\node[thm, below right=of L1] (AT) {Formule AT\\3 régimes\\{\scriptsize\THMTAG}};
\node[idstyle, below=of AT] (C1) {C1~: ph+st\\$= n_{\mathrm{base}}$\\{\scriptsize\IDTAG}};
\node[thm, below=of C1] (PMPT) {MP--PT~:\\$\mathrm{AT}=1$, $p\ge 41$\\{\scriptsize\THMTAG}};

\draw[arr] (L1) -- (dim);
\draw[arr] (L1) -- (sig);
\draw[arr] (L1) -- (AT);
\draw[arr] (C1) -- (AT);
\draw[arr] (AT) -- (PMPT);
\end{tikzpicture}
\end{center}

\noindent
\textbf{Niveaux de preuve.}\enspace
Le tableau~\ref{tab:PM_proof_status} résume le statut de preuve et
les dépendances de chaque résultat.

\begin{table}[htbp]
\centering
\caption{Résultats MP~: statut de preuve.}
\label{tab:PM_proof_status}
\begin{tabular}{lll}
\toprule
\textbf{Résultat} & \textbf{Statut} & \textbf{Dépendances} \\
\midrule
L1 (incrément unité) & \THMTAG & Algèbre des indicateurs + TCR \\
$\Sigma_k = \binom{4}{k}$ & \THMTAG\ (pos.\ gén.) & AT-position générale \\
$\dim(d) = P(d{+}1) - (2d{+}1)$ & \THMTAG\ (pos.\ gén.) & Décomp.\ par bloc \\
AT parties A, B & \THMTAG & Réemploi canal TCR + L1 \\
AT partie C & \THMTAG\ (pos.\ gén.) & C1 \IDTAG\ + C2 + canaux TCR \\
MP--PT ($\mathrm{AT} = 1$, $p \ge 41$) & \THMTAG & Corollaire de la partie~A \\
2$^{\rm e}$ instanciation (codes) & \VALTAG & 7 codes, 3 modes de cond. \\
\bottomrule
\end{tabular}
\end{table}

% ---- Chapter ledger ------------------------------------------
\begin{ledgerbox}
\begin{tabular}{lll}
Script & Tests & Score \\
\midrule
test\_L1.py & Incrément L1 sur 13 strates & 13/13 \\
test\_dim\_formula.py & $\dim(d) = P(d{+}1) - (2d{+}1)$ & 10/10 \\
test\_sigma\_spectrum.py & $\Sigma_k = \binom{4}{k}$ & 13/13 \\
test\_AT\_formula.py & Formule AT en 3 régimes & 13/13 \\
test\_phantom\_rank.py & Identité C1 + saturation C2 & 6/6 \\
test\_codes\_instance.py & 2$^{\rm e}$ instanciation (codes $\mathbb{F}_2$) & 7/7 \\
test\_pm\_lucky\_deep.py & 3$^{\rm e}$ instanciation (nombres chanceux) & 15/15 \\
\end{tabular}
\end{ledgerbox}


\newpage

\section{Extensions de convergence~: course arithmétique, graphes spectraux et symétrie}
\label{sec:convergence_ext}

\epigraph{L'art de faire des mathématiques consiste à trouver ce cas
particulier qui contient tous les germes de la généralité.}{David
Hilbert}

% ---- Status Box (R1) ----------------------------------------
\begin{statusbox}
\textbf{OBJECTIF~:}
  Présenter trois analyses étendues qui approfondissent et valident le
  théorème de convergence T4 de l'article
  compagnon~\cite{companion_M3}~: (i)~la course
  arithmético-géométrique, (ii)~la théorie spectrale des graphes du
  crible, et (iii)~le groupe de symétrie $D_4$.

\textbf{ENTRÉES~:}
  T4 (théorème~T4~\cite{companion_M3}, article
  compagnon~\cite{companion_M3}), structure TCR (article
  compagnon~\cite{companion_M1}), décomposition spectrale de $T_3$
  (article compagnon~\cite{companion_M3}).

\textbf{RÉSULTATS PRINCIPAUX~:}
  \begin{itemize}[nosep]
  \item Course arithmético-géométrique~: cinq routes alternatives vers
    T4, toutes cohérentes.
  \item $K_3$ est un expanseur de Ramanujan optimal~; sa fonction zêta
    d'Ihara a tous ses zéros non triviaux sur le cercle critique
    $|u| = 1/\sqrt{q}$.
  \item $D_4 = \langle T_3, J \rangle$ organise la structure
    spectrale~; $J$ entrelace les deux secteurs propres
    $v_+, v_-$.
  \end{itemize}

\textbf{STATUT~:}
  \THMTAG~(expanseur, $D_4$) + \VALTAG~(course, bornes spectrales)

\textbf{NON PROUVÉ~:}
  Les routes~C, D, E du paysage de preuve sont actives (contrôles de
  cohérence supplémentaires, non requis pour la clôture T4).
\end{statusbox}

% ============================================================
\subsection{La course arithmético-géométrique}
\label{sec:ch7b_race}

L'investigation a substantiellement renforcé l'argument T4 et
l'évidence de convergence. Les résultats clés sont résumés ici.

\subsubsection{Formule exacte des 2-grammes TCR (THM)}

La formule TCR de mise à jour pour les comptes de 2-grammes (théorème
de mise à jour TCR~\cite{companion_M1}) a été vérifiée exactement
jusqu'à $k = 11$, par calcul exhaustif sur
$P(10) = 6\,469\,693\,230$ résidus
($N(10) = 1\,021\,870\,080$ survivants) et propagation TCR à
$k = 11$ ($N(11) = 30\,656\,102\,400$).
La formule de facteur $(p-3)$ produit une erreur nulle dans tous les
cas.

\emph{Découverte critique}~: la formule analogue pour les comptes de
3-grammes, $n'_3(a,b,c) = (p-4)\,n_3(a,b,c) + A_3 + B_3$, est
\textbf{fausse}. Le facteur multiplicatif exact est non entier
(par ex., $3{,}0$, $7{,}75$, $9{,}54$, $\ldots$ pour le triple
$(0,1,2)$ aux niveaux successifs). Ceci révèle une \emph{hiérarchie
de type BBGKY}~: la fermeture de la récurrence des 3-grammes
nécessite des 4-grammes, qui nécessitent des 5-grammes, et ainsi de
suite. La hiérarchie ne se ferme à aucun niveau fini.

\subsubsection{Le couplage $D \approx A_{10}$ des trois-grammes}
\label{sec:ch7b_DA10}

La formule TCR de mise à jour pour les 2-grammes décompose la
correction de récurrence $\Delta(k)$ en termes de bord. On définit
le \emph{compte des trois-grammes}
\begin{equation}
  A_{10}(K) \;=\; \#\bigl\{\,(i,i{+}1,i{+}2) \in \text{survivants}
  : s_i \equiv 1,\; s_{i+1}+s_{i+2} \equiv 0 \pmod{3}\bigr\}.
  \label{eq:A10_def_ext}
\end{equation}
La correction se sépare en
$\mathrm{Total} = D + 2A_{10}$
où $D = n_{12}-n_{10}$ est le biais dominant.

\begin{proposition}[Quasi-identité $D = A_{10}$]
\label{prop:DA10_ext}
Pour toutes les profondeurs de crible calculées $K = 3,4,5$~:
\[
  D(K) \;=\; A_{10}(K) \qquad (\text{exact}).
\]
Pour $K \geq 6$ un petit résidu $D - A_{10} = O(\sqrt{N_K})$
apparaît, cohérent avec des corrections de mélange.
L'identité $D = A_{10}$ montre que le biais dominant T4 $D > 0$ est
entièrement engendré par un motif de trois-grammes spécifique~--- c'est
un compte \emph{structurel}, et non une fluctuation.
\end{proposition}

\subsubsection{La correction non-markovienne $\eta$ et le coefficient $c_\eta$}

On définit la correction non-markovienne
$\eta(k) = \beta_{\rm actuel}(k) - \beta_{\rm Markov}(k)$, où
$\beta$ est la probabilité conditionnelle d'un motif de
trois-grammes spécifique. Le coefficient normalisé
$c_\eta = |\eta|/\varepsilon$ mesure la force des corrélations de
trois-grammes relativement à l'écart
$\varepsilon = 1/2 - \alpha$.

Les données exactes $k = 3\ldots 11$ montrent~:
\begin{itemize}[nosep]
\item $c_\eta$ croît~: $0{,}28$, $0{,}36$, $0{,}39$, $0{,}42$,
  $0{,}44$, $0{,}46$, $0{,}46$
\item La croissance \emph{décélère}~: les rapports
  $1{,}093$, $1{,}069$, $1{,}031$, $1{,}009$ convergent vers~$1$
\item À chaque niveau, $c_\eta < c_\eta^{\max}$ (le seuil critique),
  avec une marge déclinant de $40\%$ à $13{,}9\%$
\end{itemize}

\subsubsection{Auto-régulation (rétroaction négative)}
\label{sec:ch7b_auto}

La valeur TCR-exacte $T_{00}(11)$, calculée depuis la mise à jour des
2-grammes, coïncide avec la prédiction du «~modèle cohérent~» (qui
inclut $\eta$ dans la mise à jour d'état) à la précision machine. Cela
confirme un mécanisme de rétroaction négative~:

\begin{enumerate}[nosep]
\item $\eta < 0$ $\Rightarrow$ $T_{00}$ croît plus lentement que la
  prédiction de Markov
\item $T_{00}$ plus bas $\Rightarrow$ $c_\eta^{\max}$ reste plus haut
  $\Rightarrow$ plus de marge pour $\eta$
\item Le système ne peut s'emballer~: il s'auto-régule
\end{enumerate}

Le bonus de $c_\eta^{\max}$ (cohérent vs.\ Markov) croît de $+7{,}5\%$
à $k = 11$ à $+138\%$ à $k = 40$.

\begin{thmbox}{Formule TCR pour $n_{00}$}
\label{thm:n00_CRT}\leanTodo{PT.Stochastic.N00CRTFormula}
Le compte bigramme diagonal $n_{00}$ satisfait la récurrence
exacte~:
\begin{equation}
  n_{00}^{(k+1)} = (p_{k+1} - 3)\,n_{00}^{(k)} + 2\,S(k),
  \label{eq:n00_CRT}
\end{equation}
où
\begin{equation}
  S(k) = n_3(0,0,0) + n_3(0,1,2) + n_3(0,2,1)
  \label{eq:S_def}
\end{equation}
est le compte des 3-grammes commençant par résidu~$0$ dont les
deuxième et troisième éléments somment à $0\!\pmod{3}$.

\begin{proof}
Lors de l'extension du crible de la profondeur~$k$ à $k{+}1$ par
suppression des multiples de $p = p_{k+1}$, la période étendue
contient $p$~copies du motif de profondeur~$k$. Chaque entier
$k$-rugueux a exactement une copie supprimée (celle
$\equiv 0\!\pmod{p}$).

\textbf{Destructions.} Chaque bigramme $(0,0)$ existant est détruit par
exactement~$3$ sites de suppression~: l'élément à sa gauche, entre
ses deux sauts, ou à sa droite. Sur les $p$~copies, $(p{-}3)$ copies
laissent chaque bigramme $(0,0)$ intact, absorbant la destruction
dans le facteur $(p{-}3)$.

\textbf{Créations.} Retirer l'élément~$y$ de
$\ldots w,x,y,z,u\ldots$ fusionne les sauts $g_L = y{-}x$ et
$g_R = z{-}y$ en $g' = z{-}x$ avec $r' = (r_L + r_R)\!\bmod{3}$. Un
nouveau bigramme $(0,0)$ est créé à \emph{gauche} lorsque
$r_{\rm FL} = 0$ et $r' = 0$ (c.-à-d.\ $r_L + r_R \equiv 0$), et à
\emph{droite} lorsque $r' = 0$ et $r_{\rm FR} = 0$. Les trois
solutions $(r_L, r_R) \in \{(0,0),\,(1,2),\,(2,1)\}$ donnent~:
\begin{align}
  \text{Créations à gauche} &= \textstyle\sum_{r_{\rm FR}}
    \bigl[n_4(0,0,0,r_{\rm FR}) + n_4(0,1,2,r_{\rm FR})
    + n_4(0,2,1,r_{\rm FR})\bigr] \notag\\
  &= n_3(0,0,0) + n_3(0,1,2) + n_3(0,2,1) = S(k).
\end{align}
Par la symétrie de renversement du temps
$n_3(a,b,c) = n_3(c,b,a)$ (vérifiée exactement pour tout
$k = 2\ldots 8$), les créations à droite valent aussi $S(k)$. Donc
$A_{00} = 2\,S(k)$.

Scripts~: \texttt{ch07\_convergence/test\_T00\_CRT\_formula.py}
(57/57 PASS).
\end{proof}
\end{thmbox}

\begin{remarque}[Statut de la règle d'évolution de $T_{00}$]
\label{rem:T00_status}
Le théorème~\ref{thm:n00_CRT} donne une récurrence exacte pour
$n_{00}$ en termes de comptes de 3-grammes. Elle \emph{ne se ferme
pas} comme fonction de $(n_{00}, \alpha, p)$ seuls, car $S(k)$ dépend
des statistiques de 3-grammes. Cependant, la formule ferme la
hiérarchie au niveau des 3-grammes~--- aucun 4-gramme n'est
nécessaire. C'est une amélioration stricte par rapport à la situation
antérieure où la mise à jour de $n_{00}$ semblait nécessiter une
information positionnelle.

La borne $T_{00} \leq \alpha$ requiert $S(k) \leq S_{\max}(k)$, où
$S_{\max} = [\alpha(k{+}1)\,n_0(k{+}1) - (p{-}3)\,n_{00}(k)]/2$. La
marge $S_{\max}/S - 1$ croît de $100\%$ à $k = 3$ à $167\%$ à
$k = 7$, confirmant l'auto-régulation.

Ce \textbf{n'est pas} une lacune logique dans la preuve de
convergence. Ce que T4 exige effectivement est~:
\begin{enumerate}[nosep]
\item $T_{00} \leq \alpha$ pour tout $k \geq 3$
  (prouvé~: induction conjointe, lemme~B)~;
\item $Q > 0$ structurellement
  (prouvé~: $b \leq n_{10}$)~;
\item $\Delta T = T_{12} - T_{00} > 0$ pour tout $k$
  (prouvé~: conséquence de~(1) et de l'identité
  $T_{12} - T_{00} = [1 - 3\alpha + 2\alpha T_{00}]/(1-\alpha)$).
\end{enumerate}
Les trois sont établis. La clôture spectrale (théorème de clôture
spectrale~\cite{companion_M3}) assure ensuite $f_{\rm bnd} \to 0$
sans jamais requérir la règle de mise à jour explicite de $T_{00}$.
Le mécanisme d'auto-régulation décrit ci-dessus fournit l'intuition
physique du \emph{pourquoi} ces bornes tiennent~: $T_{00}$
s'auto-régule par rétroaction négative, empêchant le système de
quitter le bassin de convergence.
\end{remarque}

\subsubsection{La course~: pourquoi l'écart est structurel}

Deux contractions concurrentes déterminent $c_\eta(k)$~:
\begin{itemize}[nosep]
\item \textbf{Arithmétique} (TCR)~: contracte $\eta$ au taux
  $\sim 1/p$ (lent mais infini~: $\sum 1/p$ diverge)
\item \textbf{Géométrique} (canal de Markov)~: contracte
  $\varepsilon$ au taux $C(k) \to 1$ (rapide mais s'épuise lorsque
  $\alpha \to 1/2$)
\end{itemize}
Le rapport $c_\eta = |\eta|/\varepsilon$ croît lorsque $C(k)$
(géométrique) bat la contraction arithmétique. Le
«~retournement~» se produit lorsque la contraction arithmétique
finit par dominer. La décélération
$1{,}093 \to 1{,}009$ extrapole vers un retournement près de
$k \approx 10$, après quoi $c_\eta$ décroîtrait.

\textbf{L'écart résiduel est un problème de fermeture de
hiérarchie}~: prouver le retournement requiert de borner $\eta$
algébriquement, ce qui requiert de fermer la hiérarchie des
3-grammes~--- un problème de type BBGKY structurellement analogue au
problème de fermeture en mécanique statistique.

\subsubsection{Clôture spectrale}

Le mur de la hiérarchie BBGKY est contourné par le zéro structurel
$r_2(0) = 0$~: le vecteur propre antisymétrique de $T_3$ n'a pas de
composante à l'état~$0$, donc le mode spectral dominant
$\lambda_2^2 = T_{12}^2$ est annihilé depuis les corrélations
ligne~$0$ qui gouvernent $\Delta(k)$. La fraction de bord résultante
$f_{\rm bnd} = |A|\,\lambda_1^2/\Delta T$ est vérifiée $< 1$ sur les
9~niveaux et converge vers~$0$ (95/95 tests PASS).
Voir le théorème de clôture spectrale~\cite{companion_M3} pour
l'énoncé complet.

% ============================================================
\subsubsection{Paysage de preuve~: cinq routes alternatives vers T4}
\label{sec:ch7b_five_routes}

Cinq approches alternatives pour fermer la hiérarchie de convergence
ont été explorées. Elles fournissent à la fois une compréhension
structurelle et des contrôles de cohérence indépendants pour
l'argument de clôture spectrale.

\begin{center}
\begin{tabular}{clll}
\toprule
\textbf{Route} & \textbf{Nom} & \textbf{Statut} & \textbf{Obstacle clé / résultat} \\
\midrule
A & 3-gramme TCR & \textbf{Close} & Facteur d'échelle non entier~; pas de forme affine \\
B & Borne indép.\ $\alpha$ & \textbf{Close} & Pas de borne $\alpha < 0{,}45$ sans T5~; $R < 2$ partiel \\
C & Amplif.\ $\Delta_{\rm Markov}$ & Active & $\xi = \Delta/\Delta_M$ monotone $1{,}5 \to 6{,}6$ \\
D & Induction $T_{00}$ + renversement temps & Active & $n_3(a,b,c) = n_3(c,b,a)$ exact~; $\Gamma < 0$ pour $k \geq 5$ \\
E & Condition maîtresse & Active & $D > 0 \Leftrightarrow \alpha(1-T_{00}) < (1-\alpha)/2$~; marge $19\%$ \\
\bottomrule
\end{tabular}
\end{center}

\paragraph{Route A (S15.6.283).}
Une tentative d'étendre la formule TCR des 2-grammes aux 3-grammes
échoue~: le facteur d'échelle $(p-3)$ pour les 2-grammes devient un
rapport non entier pour les 3-grammes (par ex.\ $7{,}75$ à
$p = 11$). Aucune récurrence affine n'existe au niveau des 3-grammes.

\paragraph{Route B (S15.6.284).}
Quatre identités exactes (I1--I4) reliant $Q$, $R$, $P$, $h$ sont
prouvées, donnant $Q > 4\varepsilon$ et $R < 2$ pour $k \geq 4$.
Cependant, aucune borne indépendante $\alpha < 0{,}45$ ne peut être
établie sans supposer T5 lui-même.

\paragraph{Route C (S15.6.285).}
La prédiction de Markov
$\Delta_M = n_{10}(1 - 3T_{00}) - D^2/n_1$
est systématiquement dépassée par le $\Delta$ exact. Le facteur
d'amplification non-markovien $\xi = \Delta/\Delta_M$ croît
monotonement de $1{,}53$ ($k = 4$) à $6{,}63$ ($k = 9$), entraîné par
la contrainte T0 $n_3(1,0,1) = 0$ qui amplifie les termes favorables
d'un facteur $\sim 1{,}93$.
(Script~: \texttt{test\_T4\_route\_C\_direct.py}, 15/15 PASS.)

\paragraph{Route D (S15.6.286).}
Une nouvelle symétrie structurelle est découverte~: l'identité de
renversement du temps $n_3(a,b,c) = n_3(c,b,a)$ tient
\emph{exactement} à tous les niveaux $k = 2\ldots 9$. Une induction
sur $E = n_{01} - n_{00}$ donne
$E(k+1) = (p-3)E(k) + \Gamma(k)$,
mais $\Gamma < 0$ pour $k \geq 5$ empêche une clôture directe.

\paragraph{Route E (S15.6.287).}
La condition maîtresse
$D > 0 \Leftrightarrow \alpha(1-T_{00}) < (1-\alpha)/2$
définit une courbe d'équilibre $\alpha_{\rm eq} = 1/(3-2T_{00})$. Les
$\alpha_k$ réels restent $19\%$--$50\%$ en dessous de
$\alpha_{\rm eq}$ pour $k = 3\ldots 9$, montrant que le système est
bien à l'intérieur du bassin de convergence.

\medskip\noindent
La clôture spectrale (mécanisme de la route~C combiné au théorème de
clôture spectrale~\cite{companion_M3}) est la route qui ferme T4
formellement. Les routes~D et~E fournissent une validation
structurelle indépendante.

% ============================================================
\subsubsection{Preuve alternative~: borne spectrale sur $S(k)$}
\label{sec:ch7b_S_spectral_bound}
% S15.6.316

La route~C peut être reformulée en une preuve TCR \emph{directe} de
$T_{00} \leq \alpha$, court-circuitant entièrement la machinerie de
$f_{\rm bnd}$. La chaîne est~:
\[
  S(k) < S_{\max}(k)
  \quad\Longrightarrow\quad
  n_{00}^{(k+1)} \leq \alpha_{k+1}\,n_0^{(k+1)}
  \quad\Longrightarrow\quad
  T_{00}^{(k+1)} \leq \alpha_{k+1}.
\]

\begin{thmbox}{Borne spectrale sur $S$}
\label{thm:S_spectral_bound}\leanProved{PT.Stochastic.SpectralDominance.T3\_spectral\_radius\_eq\_one}
Pour tout $k \geq 3$, $S(k) < S_{\max}(k)$, où
$S_{\max} = [\alpha_{k+1}\,n_0^{(k+1)} - (p_{k+1}-3)\,n_{00}^{(k)}]/2$.

\begin{proof}
\textbf{Étape~1~: décomposition de Markov.}
On écrit $S = S_{\rm M} + \delta S$, où l'approximation de Markov
est
\[
  S_{\rm M} = n_0\bigl(T_{00}^2 + 2\,T_{01}\,T_{12}\bigr)
\]
et
$\delta S = \sum_{(b,c)\in\mathcal{S}} [n_3(0,b,c) - n_0\,T_{0b}\,T_{bc}]$
avec $\mathcal{S} = \{(0,0),\,(1,2),\,(2,1)\}$.

\textbf{Étape~2~: contrôle spectral.}
L'écart $|\delta S|/n_0 \leq K_S\,|\lambda_1|^2$ avec
$K_S \leq 4{,}7$ (borne effective, vérifiée $k = 3\ldots 9$). Cela
découle de la structure spectrale de~$T_3$~: le seul mode propre
contribuant aux corrélations à deux pas depuis l'état~$0$ est
$\lambda_1 = (T_{00}-\alpha)/(1-\alpha) \to 0$, le mode
antisymétrique $\lambda_2 = -T_{12}$ étant annihilé par
$r_2(0) = 0$.

\textbf{Étape~3~: domination de la marge.}
La marge de Markov
$\mu = (S_{\max} - S_{\rm M})/n_0 = O(\varepsilon)$
domine la correction $|\delta S|/n_0 = O(\varepsilon^2)$
(puisque $|\lambda_1| = O(x/(1{-}\alpha))$ et
$x/\varepsilon < 1$ est monotone décroissant pour $k \geq 5$).

Numériquement, $\mu/|\delta S/n_0| \geq 2{,}1$ à $k = 3$ et croît
à $\geq 11{,}8$ à $k = 8$, confirmant l'asymptotique
$\mu/|\delta S| \sim 1/\varepsilon \to \infty$.

Donc $S = S_{\rm M} + \delta S < S_{\rm M} + \mu\,n_0 = S_{\max}$.

Scripts~: \texttt{ch07\_convergence/test\_T00\_S\_spectral\_bound.py}
(46/46 PASS).
\end{proof}
\end{thmbox}

\begin{remarque}[Indépendance vis-à-vis du lemme de mélange]
Cette preuve partage les mêmes ingrédients spectraux (annihilation
$r_2(0) = 0$, dilution TCR, décroissance des valeurs propres) que le
lemme de mélange (théorème de clôture
spectrale~\cite{companion_M3}), mais la chaîne logique
$S < S_{\max} \Rightarrow T_{00} \leq \alpha$ est indépendante de
l'argument $f_{\rm bnd} < 1$. Elle fournit une seconde clôture,
structurellement distincte, de la borne $T_{00}$.
\end{remarque}

% ============================================================
\subsection{Théorie spectrale des graphes du crible}
\label{sec:ch7b_spectral_graph}

Les arguments spectraux des sections précédentes peuvent recevoir une
fondation \emph{théorique des graphes} qui rend le mécanisme de
contraction transparent.

\subsubsection{Le graphe du crible comme expanseur optimal}

Au module $m = 3$, les transitions permises forment le graphe complet
$K_3$ sur $\{0,1,2\}$ (chaque sommet relie aux deux autres, puisque
$T_{11} = T_{22} = 0$ ne supprime que les boucles).
Soit $L = D - A$ son laplacien de graphe, où $D$ est la matrice des
degrés et $A$ la matrice d'adjacence.

\begin{thmbox}{$K_3$ est un expanseur optimal}
\label{thm:expander_ext}\leanTodo{PT.Graph.K3OptimalExpander}
Le graphe de transition modulo~3 a~:
\begin{enumerate}[nosep]
\item Spectre laplacien $\operatorname{spec}(L) = \{0, 3, 3\}$,
  donnant $\lambda_2(L) = 3$ (le maximum possible pour un graphe à
  3~sommets).
\item Constante de Cheeger $h = 1$ (le maximum pour tout graphe)~:
  toute coupure de sommet sépare le graphe en parties à peu près
  égales.
\item Le rapport isopérimétrique $h_{\min} \geq 1/2$ est garanti pour
  tout $K \geq 3$ par les contraintes structurelles
  $T_{11} = T_{22} = 0$ et
  $T_{12} \geq (p-3)/(p-1) \geq 1/2$.
\end{enumerate}
\end{thmbox}

\begin{proof}
$K_3$ est $2$-régulier ($d_{\max} = 2$), donc $D = 2I$.
La matrice d'adjacence a $a_{ij} = 1$ pour $i \neq j$, donc
$A = J - I$ où $J$ est la matrice toute-à-un.
Le laplacien est $L = D - A = 2I - (J - I) = 3I - J$.
Les valeurs propres de $J$ sur $\mathbb{R}^3$ sont $\{3, 0, 0\}$,
donnant $\operatorname{spec}(L) = \{0, 3, 3\}$.
Puisque $\lambda_2(L) = 3 > d_{\max} = 2$, $K_3$ sature la borne
serrée $\lambda_2 \leq 2d_{\max} = 4$ pour les laplaciens de graphes.
La constante de Cheeger est $h = 1$, atteinte par toute coupure à un
sommet $S = \{v\}$~: $|\partial S| = 2$ arêtes quittent $S$, et
$h = |\partial S|/|S| = 2/2 = 1$.
L'inégalité de Cheeger pour les graphes $d$-réguliers,
$\lambda_2/(2d) \leq h \leq \sqrt{2\lambda_2/d}$, donne
$3/4 \leq h \leq \sqrt{3}$, cohérent avec $h = 1$.
Scripts~: \texttt{ch\_math\_structures/test\_graph\_laplacian.py}
(31/31 PASS).
\end{proof}

\begin{remarque}[Pourquoi la propriété d'expanseur compte]
\label{rem:expander_matters_ext}
La constante de Cheeger $h > 0$ \emph{uniforme en $K$} implique que
$|\lambda_2(T_K)| \leq 1/(1 + h_{\min}) < 1$ pour tous les niveaux du
crible. C'est une route indépendante vers la contraction
$|\lambda_2| < 1$ qui ne repose pas sur le calcul explicite des
valeurs propres à chaque niveau. Elle découle des contraintes
structurelles C2 ($T_{11} = T_{22} = 0$) et C4 ($T_{12} > 0$).
\end{remarque}

\subsubsection{Fonction zêta d'Ihara}
\label{subsec:ch7b_ihara}

La fonction zêta d'Ihara du graphe du crible fournit un encodage
multiplicatif de sa structure cyclique~:
\begin{equation}
  Z_G(u) = \prod_{[C]} (1 - u^{|C|})^{-1},
  \label{eq:ihara_ext}
\end{equation}
où le produit s'étend aux cycles primitifs~$[C]$.
Pour $K_3$, la formule de Bass--Hashimoto donne une expression
rationnelle explicite dont les zéros sont les racines cubiques de
l'unité. Tous les zéros non triviaux se trouvent sur le cercle
critique $|u| = 1/\sqrt{q}$, confirmant que $K_3$ est un expanseur
de Ramanujan optimal.

Scripts~: \texttt{ch\_math\_structures/test\_ihara\_zeta.py}
(16/16 PASS).

\subsubsection{Borne spectrale uniforme~: trois lacunes fermées}
\label{subsec:ch7b_three_gaps}

La preuve formelle de $|\lambda_2(T_K)| < 1$ pour tout $K$ requiert
de fermer trois lacunes potentielles~:

\begin{enumerate}[leftmargin=*]
\item \textbf{Lacune~1 (uniformité)~:} a-t-on $h_{\min} > 0$ pour
  tout~$K$~? \textbf{FERMÉE.} La contrainte structurelle
  $T_{12}(K) \geq (p_K - 3)/(p_K - 1) \geq 1/2$ pour $K \geq 3$
  garantit une constante de Cheeger positive à chaque niveau.

\item \textbf{Lacune~2 (survivants $\to$ entiers)~:} la
  classification modulo~3 s'applique-t-elle à tous les entiers, et
  pas seulement aux survivants~?
  \textbf{DISSOUTE.} La matrice de transfert $T_K$ est une matrice
  $3 \times 3$ opérant sur les classes $\{0, 1, 2\}$. Chaque entier
  possède un résidu modulo~$p$ bien défini. Le crible classifie
  \emph{tous} les entiers via leurs résidus~; la classe~$0$ (les
  multiples) est déjà incluse. Il n'y a pas de pont à construire.

\item \textbf{Lacune~3 (constante $C$)~:} la constante de corrélation
  $C$ dans $|A| \leq C$ est-elle uniformément bornée~?
  \textbf{FERMÉE.} La décomposition spectrale de la matrice de
  transfert jointe donne
  $C \leq \dim(\text{espace propre joint}) = 4$, vérifiée
  exactement comme $C \in [2{,}58~;\; 4{,}50]$ pour
  $K = 3 \ldots 11$.
\end{enumerate}

Ces trois clôtures, combinées au théorème de Perron--Frobenius et à
l'inégalité de Cheeger, donnent la borne spectrale uniforme.

Scripts~: \texttt{ch\_math\_structures/test\_spectral\_bound.py}
(15/15 PASS).

\subsubsection{Décroissance de mémoire de Liouville}
\label{subsec:ch7b_memory}

Le caractère hybride $\lambda(n) \cdot \chi_3(n)$ définit une matrice
de transfert jointe $4 \times 4$. L'opérateur de torsion $D_{01}$
(qui couple les secteurs Liouville et caractère) a un rayon spectral
$\rho(D_{01}) = 0{,}05 \ll 1$. Par conséquent, l'indice d'obstruction
\begin{equation}
  I(K) = \rho(D_{01},\, T_{\rm joint}(K)) \sim e^{-0{,}71\,K}
  \label{eq:obstruction_decay_ext}
\end{equation}
décroît géométriquement, ce qui signifie que la mémoire de la
fonction de Liouville pour ses corrélations initiales avec $\chi_3$
s'évanouit exponentiellement avec la profondeur du crible.

L'obstruction à RH identifiée dans l'article
compagnon~\cite{companion_M1} (la croissance de Liouville
$r_K \sim C \cdot K$) est \emph{localisée dans le vecteur propre
stationnaire $v_+$}, et non dans le mode contractant $v_-$. Cette
localisation explique pourquoi l'obstruction persiste en amplitude
($r_K$ croît) tandis que son contenu en information décroît
($I(K) \to 0$)~: c'est un artefact de projection, et non une
instabilité dynamique.

Scripts~:
\texttt{ch\_math\_structures/test\_liouville\_spectral.py} (12/12),
\texttt{test\_hybrid\_character.py} (10/10),
\texttt{test\_obstruction\_index.py} (10/10).

% ============================================================
\subsection{Le groupe de symétrie $D_4$ du crible}
\label{sec:ch7b_d4}

La structure spectrale du crible a un groupe de symétrie qui organise
l'argument de convergence.

\begin{thmbox}{$D_4$ comme groupe de symétrie du crible}
\label{thm:d4_ext}\leanProved{PT.Sieve.KleinFourEmbedding}
On définit l'entrelacement $J\colon f \mapsto f \cdot \chi_3$,
restreint aux résidus non nuls $\{1,2\}$ (l'espace d'états du crible).
Sur ce sous-espace de dimension~$2$, $J = \mathrm{diag}(1,-1)$ est une
involution ($J^2 = I$), et~:
\begin{enumerate}[nosep]
\item $J$ échange les vecteurs propres spectraux~:
  $J(v_+) = v_-$, $\;J(v_-) = v_+$.
\item Le groupe engendré par $T_3$ et $J$ est le groupe diédral~:
  $\langle T_3, J \rangle \cong D_4$.
\item Sous $D_4$, l'espace $\mathbb{R}^3$ des fonctions du crible se
  décompose comme $\rho_1 \oplus \rho_5$ (les représentations
  irréductibles triviale et fidèle de dimension~2).
\item $v_+$ et $v_-$ forment la base canonique de $\rho_5$.
\end{enumerate}
\end{thmbox}

\begin{proof}
L'entrelacement $J$ agit par multiplication ponctuelle avec
$\chi_3$. Sur $\{0,1,2\}$, $\chi_3$ a pour valeurs $(0, 1, -1)$~;
restreint à l'espace d'états du crible $\{1,2\}$,
$J = \mathrm{diag}(1,-1)$. Sur la base propre de $T_3$~:
$v_+ = (1, 1)/\sqrt{2}$ (symétrique, valeur propre $+1$) et
$v_- = (1, -1)/\sqrt{2}$ (antisymétrique, valeur propre $-1$).
L'action $J(v_+)(r) = \chi_3(r) \cdot v_+(r) = v_-(r)$ et
réciproquement. Puisque $T_3^2 = I$ et $J^2 = I$ avec
$T_3 J \neq J T_3$, le groupe engendré est d'ordre~$8$ et isomorphe
à $D_4$ (vérifié par table de multiplication explicite). La
décomposition irréductible découle de la table des caractères de
$D_4$ appliquée à la représentation naturelle sur $\mathbb{R}^3$.
Scripts~: \texttt{ch\_math\_structures/test\_intertwiner.py}
(17/17), \texttt{test\_d4\_representations.py} (35/35).
\end{proof}

\begin{remarque}[Symétrie d'entrelacement]
L'entrelacement $J$ échange $v_+ \leftrightarrow v_-$~; puisque
$\langle T_3, J \rangle \cong D_4$, toute affirmation spectrale sur
$T_3$ dans un secteur propre est automatiquement traduite, via $J$,
en une affirmation dans l'autre secteur.
\end{remarque}

\newpage

\section{Conclusion}
\label{sec:conclusion}

Cet article a développé l'infrastructure mathématique étendue de la
théorie de la persistance à travers quatre domaines~: structures
algébriques, mécanique complexe, mathématiques prédictives et
extensions de convergence. Nous résumons les principaux résultats et
leur statut.

\textbf{Algèbre du crible et transformées.}
Le produit du crible $*_T$ définit un objet algébrique véritablement
nouveau~--- non associatif, non commutatif, sans élément neutre~---
distinct de toutes les algèbres connues de fonctions arithmétiques.
La métrique PT $d_{\rm PT}$ est quasi-orthogonale à la fois aux
métriques archimédienne et $p$-adique, et la catégorie du crible
$\mathrm{Cat}(\mathrm{Sieve})$ admet quatre foncteurs cohérents.
Trois nouvelles transformées étendent la transformée de persistance~:
la transformée tensorielle $\mathcal{T}$ capture les corrélations
inter-modulaires dans un tenseur à $3{\times}5{\times}7 = 105$
composantes~; la transformée holonomique $\mathcal{H}$ dote l'espace
des fonctions arithmétiques d'une métrique de Fisher et d'un produit
eulérien généralisé~; et la transformée de décohérence $\mathcal{D}$
établit le théorème~H, seconde loi de la thermodynamique pour le
crible (l'entropie de la distribution des classes de sauts est non
décroissante). Le code des résidus de sauts atteint la distance de
Hamming $d = n - 1$ (asymptotiquement MDS), identifiant prédiction et
correction d'erreur~: deux faces de la même structure mathématique.

\textbf{Mécanique complexe.}
Le relèvement à la variable complexe $w_p$ révèle que l'espace
d'états du crible est un cercle de rayon $s = 1/2$~--- le même
paramètre qui gouverne la distribution stationnaire. L'identité
$|w_p|^2 = \mathrm{Re}(w_p)$ relie la règle de Born à la géométrie
circulaire. La force holomorphe, l'intégrabilité de Liouville, la
relation d'incertitude exacte et la décomposition des corrections
radiatives QED en quantités PT-natives ($N_{\mathrm{spatial}} = 3$,
$s = 1/2$, $\mathrm{Circ}(\mathcal{C}) = \pi$) montrent que le plan
complexe est l'habitat naturel de la dynamique du crible.

\textbf{Mathématiques prédictives~: sept théorèmes.}
Le cadre MP fournit un langage universel pour les cascades de
contraintes, instancié sur trois domaines~:
\begin{enumerate}[nosep]
\item \textbf{L1} (incrément unité)~: chaque premier ajoute
  exactement un degré de liberté. Inconditionnel~; prouvé à partir
  de l'algèbre des indicateurs et de l'irréductibilité TCR.
\item \textbf{Formule de dimension}~: $\dim(d) = P(d{+}1) - (2d{+}1)$,
  vérifiée sur 10~couches, prédite pour la 11$^{\rm e}$.
\item \textbf{Spectre d'activation}~: $\Sigma_k = \binom{4}{k}$ pour
  $k \geq \mathrm{AT}$, vérifié sur 13~strates.
\item \textbf{Formule AT} (trois régimes)~: écho ($|B| \leq 2$,
  $\mathrm{AT} = 1$), construction ($3 \leq |B| \leq 8$,
  $\mathrm{AT} = \lceil(|B|-2)/3\rceil + 1$), asymptotique
  ($|B| \geq 9$, $\mathrm{AT} = 1$).
\item \textbf{Identité de partition C1}~: $\mathrm{fantome}(1)
  + \mathrm{struct}(1) = n_{\mathrm{base}}$, exacte sur 6~strates.
\item \textbf{Saturation fantôme C2}~: pour $|B| \geq 3$, le fantôme
  absorbe tous les termes croisés.
\item \textbf{Corollaire MP--PT}~: $\mathrm{AT} = 1$ pour tout
  $p \geq 41$, prouvant que la complexité structurelle est épuisée
  dans le premier cycle primoriel.
\end{enumerate}
La constante $N_{\mathrm{spatial}} = 3$ apparaissant dans la formule
AT est la même quantité qui détermine la dimensionnalité spatiale
via le critère des premiers actifs~\cite{companion_M2}~: c'est la
constante d'obstruction à la détection structurelle à travers le
bruit arithmétique.

\textbf{Extensions de convergence.}
La course arithmético-géométrique révèle cinq routes alternatives
vers T4, toutes cohérentes. Le graphe de transition modulo~3 $K_3$
est un expanseur de Ramanujan optimal ($\lambda_2(L) = 3$, constante
de Cheeger $h = 1$), et sa fonction zêta d'Ihara place tous ses
zéros non triviaux sur le cercle critique $|u| = 1/\sqrt{q}$. Le
groupe diédral $D_4 = \langle T_3, J \rangle$ organise la structure
spectrale, avec l'entrelacement $J$ échangeant les deux secteurs
propres $v_+$ et $v_-$.

\textbf{Interprétation théorique des codes.}
Le crible est simultanément un algorithme générateur de premiers, un
canal de décohérence CPTP, un code correcteur d'erreurs de distance
$d = n - 1$, et une contraction spectrale de trou $1 - |\lambda_2|$.
La distance de code $d = n - 1$ est la manifestation arithmétique de
la persistance informationnelle~: le crible produit un code dont la
redondance croît sans borne, garantissant que le motif des premiers
est récupérable à partir d'une information partielle. Le code
quantique TCR $[[3, 5.58, 1]]$ d'écart de Knill--Laflamme
$\Delta_{\rm KL} = 0{,}008$ montre que prédiction et correction
d'erreur sont la même structure mathématique vue sous des angles
différents.

\textbf{Question ouverte.}
La conjecture de position générale (conjecture~\ref{hyp:gen_pos})~---
selon laquelle aucune relation linéaire accidentelle ne survient
parmi les termes croisés TCR de la matrice d'observation~--- est
vérifiée numériquement pour les 13~strates explorées mais n'est pas
dérivée à partir de premiers principes arithmétiques. Une preuve
promouvrait la formule de dimension et le spectre d'activation, du
statut conditionnel à celui de théorèmes inconditionnels.

\textbf{Suite.}
L'article compagnon~\cite{companion_M5} construit le pont de
l'arithmétique vers la physique~: la dualité de Pontryagin sur la
décomposition TCR, la contraction de Buchstab et la reconstruction
d'Osterwalder--Schrader fournissent un chemin formel des structures
du crible développées ici vers le couplage produit
$\alpha_{\mathrm{sieve}} = \prod_{p \in \{3,5,7\}} \sinsp{p}$ et son
interprétation physique.

\bibliographystyle{plainnat}
\bibliography{references}

\end{document}
